\begin{comment}
\documentclass[letter,11pt]{amsart}
\usepackage{verbatim}
\usepackage{amsmath}
\usepackage{amssymb, latexsym}
\usepackage{amsthm}
\usepackage{amscd}
\usepackage{amsfonts,mathrsfs}
\usepackage{delimseasy1.01}
\setcounter{tocdepth}{3}

\newcounter{keepenumi}%doesn't belong here, but where??
\newcommand{\enumiget}{\setcounter{enumi}{\value{keepenumi}}}
\newcommand{\enumisave}{\setcounter{keepenumi}{\value{enumi}}}%
\newcounter{keepenumii}%doesn't belong here, but where??
\newcommand{\enumiiget}{\setcounter{enumii}{\value{keepenumii}}}
\newcommand{\enumiisave}{\setcounter{keepenumii}{\value{enumii}}}%
\newcommand{\pii}{\pi i}
\newcommand{\N}{\text{$\mathbb N$}}
\newcommand{\Z}{\text{$\mathbb Z$}}

\newtheorem{theorem}{Theorem}%[part]
\newtheorem{conj}[theorem]{Conjecture}
\newtheorem{ques}[theorem]{Question}
\newtheorem{corollary}[theorem]{Corollary}
\newtheorem{lemma}[theorem]{Lemma}
\newtheorem{proposition}[theorem]{Proposition}
\numberwithin{theorem}{section}
%%%%%%%%%
%Supply  abbreves.
%No problems with numbering as long as
%  EVERYTHING
% uses one sequence
%\theoremstyle{plain}
\newtheorem{thm}[theorem]{Theorem}
\newtheorem{cor}[theorem]{Corollary}
\newtheorem{lem}[theorem]{Lemma}

\newtheorem{prop}[theorem]{Proposition}
%
\newcommand\newnumbered{\newtheorem}
% remark style theorems
% remark style theorems7
% definition style theorems
\theoremstyle{definition}
%
\newtheorem{definition}[theorem]{Definition}
%\numberwithin{definition}{section}
\theoremstyle{remark}
\newtheorem{notation}[theorem]{Notation}
\newtheorem{example}[theorem]{Example}
\newtheorem{remark}[theorem]{Remark}
\newtheorem{rem}[theorem]{Remark}
\newtheorem{remarks}[theorem]{Remarks}
\newtheorem{rems}[theorem]{Remarks}
% equations use (sub)section as a parent
\numberwithin{equation}{section}
\newcommand\aset[1]{\newline$\{#1\}$}
\newcommand{\thmref}[1]{Theorem~\ref{#1}}
\newcommand\refthm\thmref
\newcommand{\secref}[1]{\S\ref{#1}}
\newcommand{\refeqn}{\eqref}
\newcommand{\conjref}[1]{Conjecture~\ref{#1}}
%
\newcommand{\refitem}[1]{(\ref{#1})} 
\newcommand{\itemref}{\refitem}
\newcommand\itref\itemref
%defined here so can change if docclass changes item numbering in enumerations
%
\newcommand{\lemref}[1]{Lemma~\ref{#1}}
\newcommand{\corref}[1]{Corollary~\ref{#1}}
\newcommand{\propref}[1]{Proposition~\ref{#1}}
\newcommand{\defref}[1]{Definition~\ref{#1}}
\newcommand{\remref}[1]{Remark \ref{#1}}
\newcommand{\remsref}[1]{Remarks \ref{#1}}
%
\newcommand\ex[2]{e^{#1\pi i/#2}}
\newcommand \One[1]{{\pmb 1}_{#1}}
\newcommand\supp{\text{Supp\,}}
\newcommand\Supp{\text{Supp\,}}
\newcommand\Times{\text{$\times$}} % lazy me
%\newcommand\Z[1]{\texttt{$\mathbb Z_2^{\,#1}$}}
\newcommand\Zp[2]{\texttt{$\mathbb Z_{#1}^{\,#2}$}}


\title[Appendix to ``A beastiary\dots'' -- DRAFT -- \today]{Appendix to ``A beastiary of sets having extremal Sidon constant, or,  there must be more
than one theorem somewhere here'' \\ \  \\ DRAFT -- \today --DRAFT }
\author{Colin C. Graham -- \today}
\address{\#208 -- 6100 Sixth Ave., Whitehorse YT Y1A 1M5 Canada}
\thanks{  
  }
\begin{document}
\begin{abstract}
This appendix contains  details not intended for publication in  
the paper  of which
it is an appendix,
but which details might be useful in the public record.
%A lengthy disclaimer concludes the appendix.
\end{abstract}

\maketitle

\centerline{\today}
\tableofcontents

%%%%%%%%%%%%%%%%%%%%%%%%%%%%%%%%%%%%%%%%%%%%%%%
%\break
\end{comment}
\addtocontents{toc} {\protect\vspace{10pt} }
\addcontentsline{toc}{section}{\qquad\qquad{\textsc{Appendix}}}
\addtocontents{toc}{\protect\vspace{2pt}}

\subsection*{Note to the referee and other readers} This appendix contains material that is not 
intended for publication because it has too many details, is too reptitive
 and also duplicative of material in the main paper. It is planned to be 
included with the pdf electronicly distributed by the author 
so that readers can see the results of the computer 
programs in fuller detain than can be justified in a publication, even an electronic one.

\section{Proofs of various items}
Here is the special case of 
\cite[Thm. 4.1]{MR2363058} that we need.

\begin{prop}\label{propAutosOfZphatk} Let $1\le m$ and $p$ a prime. then
$L=\Z_{p^m}$ has $p-1$ automorphisms.
\end{prop}

\begin{proof}
Every element $x\in L$ has the form
\[
x=
a_0+a_1p+\cdots a_{m-1}p^{m-1}
\]
where $0\le a_k<p$ for $0\le k <m$.
Every automorphism $T:L\to L$ is determined by $T(1)$.  Using the facts that $1$ is
a generator of $L$ and $T$ is an automorphism, we see that 
\begin{align}
T(x) &= T\prn{a_0+a_1p+\cdots+a_{m-1}p^{m-1}}
\\&
= T(a_0)+T(a_1)p+\cdots+T(a_{m-1})p^{m-1}.
\end{align}
\end{proof}

\begin{prop}There is no real extreme measure on $\{0,1,2,4\}\subset\Zp 7{}.$
\end{prop}
\begin{proof}
\textit{Look at the 4 possible measures.}
Case I.

$\mu=\delta_0+\delta_1+\delta_2-\delta_4.$
$M=\begin{pmatrix}
1&1&1&0&-1&0&0\\
0&1&1&1&0&-1&0\\
\dots
\end{pmatrix}
$

Row 2 inner row 1 is  $=2\ne0$.

\smallskip

Case II.

$\mu=\delta_0+\delta_1-\delta_2+\delta_4.$
$M=\begin{pmatrix}
1&1&-1&0&1&0&0\\
0&1&1&-1&0&1&0\\
0&0&1&1&-1&0&1\\
\dots
\end{pmatrix}
$

Row 3 inner row 1 is $-2=\ne 0$.

\smallskip

Case III.

$\mu=\delta_0-\delta_1+\delta_2+\delta_4.$
$M=\begin{pmatrix}
1&-1&1&0&1&0&0\\
0&1&-1&+1&0&1&0\\
0&0&1&-1&1&0&1\\
\dots
\end{pmatrix}
$

Row 2 inner row 1 is $=-2\ne 0$.

\smallskip

Case IV.


$\mu=-\delta_0+\delta_1+\delta_2+\delta_4.$

$M=\begin{pmatrix}
-1&1&1&0&1&0&0\\
0&-1&1&1&0&1&0\\
0&0&-1&1&1&0&1\\
1&0&0&-1&1&1&0\\
%1&0&1&0&0&-1&1\\
%1&1&0&1&0&0&-1\\
\dots
\end{pmatrix}
$

Row 3 inner row 1 is $=-1\ne 0$.
\qedhere
\end{proof}

\section{Proofs of extremalities for cyclic groups}
%%%%%%%%%%%%%%%%%%%%%%%%%%%%%%%%%%%%%%%%%%%%%%%
In this section we give   proofs of some of the extremalities claimed in and of related results.
 We also show that the extreme sets here are neither sums of other extreme sets nor of a subgroup and a set.

To save the reader from flipping between \cite{ABeastiary} and this document, we sometimes include details
from \cite{ABeastiary}.

\subsection{Sets with three elements}

\begin{prop}\label{prop3inZ3}
$\Zp{3}{}$ is extreme. \newline \end{prop} 


 
\begin{proof}
Let  $\mu =  \delta_{0}+e^{4\pii/3}\delta_{1}+\delta_{2}
$.  Then 
\begin{align*} 
\mu * \tilde \mu & = \big(1+1+1\big)\delta_{0}
\\&\qquad
+\big(1+e^{4\pii/3}+e^{2\pii/3}\big)\delta_{1}
\\&\qquad
+\big(e^{2\pii/3}+e^{4\pii/3}+1\big)\delta_{2}\\&=3\delta_{0}\qedhere
\end{align*}
\end{proof}



In \Zp{4}{} there is only one, up to translation, and it is extreme:
\begin{prop}\cite{MR627683}\label{prop3in4}
 $\{0,1,2\}\subset \Zp4{}$ is extreme.
\end{prop}
\begin{proof}
Let $\nu = \delta(0) + e^{3\pi i/4}\delta(1) +i\delta(2)$.
Then
\begin{align*}
\widehat\nu(0) &= 1- \frac{\sqrt2}2 + \frac{\sqrt2}2i +i, \text{ so }
\\  
|\widehat\nu(0)| &= \big|(1-\frac{\sqrt2}2)^2 +(1+\frac{\sqrt2}2)^2\big|^{1/2} =\sqrt3.
\\
\widehat\nu(1) &= 1 -\frac{\sqrt2}2 -\frac{\sqrt2}2i -i, \text{ so }
|\widehat\nu(1)|=  \sqrt3\\
\widehat\nu(2) & = 1 +\frac{\sqrt2}2 -\frac{\sqrt2}2i +i \text{, so } |\widehat\nu(2)|= \sqrt3.
\\
\widehat\nu(3) & = 1  +\frac{\sqrt2}2 +\frac{\sqrt2}2i -i, \text{ so } 
|\widehat\nu(3)| = \sqrt3.\qedhere
\end{align*}
\end{proof}

Another extreme measure is $\mu=\delta(0) - e^{3\pi i/4}\delta(1) + i\delta(2)
=\delta(0) + e^{7\pi i/4}\delta(1) + i\delta(2)$.
All extremal measures have one of the forms, $\nu$ or $\mu$:
Suppose $\mu=\delta_0+a\delta_1+b\delta_2$ is extreme on $\{0,1,2\}\subset \Zp4{}$. Then
$\mu*\tilde\mu=3\delta_0+(a+\bar ab)\delta_1 + (b+\bar b)\delta_2 +
(\bar a+a\bar b)\delta_3=3\delta_0$.
Hence, $b+\bar b=0$ so $b=\pm i$. Assume $b=i$. Then 
$\bar a-ia=0$ so\footnote{\ 
Let $a= x+iy$. Then $\bar{a}-ia= x-iy-ix+y=0$ means $x+y=0$. Since 
$|a|=1$, $x= \pm\sqrt2/2$.}
either $a=\pm e^{\pm\pi i/4}$ or $a =\pm e^{\pm 3\pi i/4}$. 
Since $y=-x$, we have $a=\pm \exp(3\pi /4).$


\medskip
Every three-element extreme set can be obtained from \Zp{3}{} and a 3-three element subset of \Zp{4}{}
by the operations of group automorphism, passing to a subgroup, and translation
\cite[3.1(ii)]{MR627683}. Here is a proof of part of \cite[3.1(ii)]{MR627683}. 
The complications of the proof here illustrate
 why the proofs
of \cite[3.1-3.3]{MR627683} occupied 200 pages of manuscript.

  
\begin{prop}[\cite{MR627683}]\label{prop3eltsets} 
If  $3\le k$ and $E=\{0,a,b\}$ is extreme in $\Zp{k}{}$ then either $E$ is a subgroup or it is a three-element
subset of a four-element subgroup.
\end{prop}

\begin{proof} Suppose $E=\{0, a,b\}\subset \Zp{k}{}$ is extreme. Let $\mu = \delta_0 +\delta_a+\delta_b.$
We may assume $0<a<b<k$. Then 

\begin{equation}\label{eq1prop3eltsets}
\mu*\tilde\mu= 3\delta_0+\delta_a+\delta_b +\delta_{-a}+\delta_{-b}+\delta_{a-b}+
\delta_{b-a}.
\end{equation}
 If $a\ne -b$ and $a\ne -a$,  then we have four pointmasses, $\delta_{\pm a},$ $ \delta_{\pm b}$ which cannot all
be matched by $\delta_{a-b},$ $\delta_{b-a}$ and $E$ is not extreme by \cite[Cor. 2.3]{ABeastiary}. 

\smallbreak
Case I: Suppose $a=-a$, that is, $a=k-a.$ Then $k=2a$ and $b\ne -b$ so $a-b\ne a+b$.
Thus,\[
\mu*\tilde\mu= 3\delta_0+2\delta_a+\delta_b +\delta_{-b}+\delta_{a-b}+
\delta_{b+a}.
\]
If $a-b=b$ then $a=2b$ and $k=4b$ so we have three elements of a  4 element subgroup, $\{0, b, 2b\}\subset\Zp{4b}{}.$
Therefore, we may assume $a-b\ne b$. 

Clearly $a-b$ is distinct from $0$, $a$, $b$, and $a+b$. Hence $\mu$ is not extreme in this case.

\smallbreak
Case II: $b=-b$, so $a\ne-a$. Then 
\[
\mu*\tilde\mu= 3\delta_0+2\delta_b +\delta_a+\delta_{-a}+\delta_{a+b}+
\delta_{b-a}
\]
or
\[
\mu*\tilde\mu= 3\delta_0+2\delta_b +\delta_a+\delta_{-a}+\delta_{a+b}+
\delta_{-b-a}.
\]
If $a=-b-a$, then $b=2a$ and $k=4a$ and we are again in the situation of a 3-element subset of a 4-element subgroup.
Therefore we may assume $a\ne -b-a$, in which case $\delta_{-b-a}$ appears only once in \eqref{eq1prop3eltsets} and $E$ is not extreme.
\smallbreak

Case III: 
$a=-b$ so 
\[ 
\mu*\tilde \mu = 3\delta_0+2\delta_a+2\delta_b + \delta_{a-b}+\delta_{b-a}.
\] 

We have two subcases. First, suppose $a= b-a$. Then since $a=-b$, $a= b--b=2b=-2a$ so $3a=0$. Hence $E=\{0,a,2a\}$ and $k=3a$, that is, $E$ is a subgroup.

Second subcase: $a-b=b-a$. We still assume $a=-b$.
  Since $a\ne b$ we must have $a-b=\ell$ has order 2, and so $k=2\ell.$ Now $a-b = \ell = -2b = 2a$ $\mod\,2\ell$. Hence, $2a=\ell,$
 and $a$ has order 4. Thus, $E=\{0,a, 2a\}\subset\Zp{4a}{}$ and $E$ is a three element subset of a 4-element subgroup. That takes care
 of the third and final case.
\end{proof}


\subsection{Sets with four elements} 

\begin{prop}\label{prop4in4}
$\Zp{4}{}$ is extreme. \newline \end{prop} 


\begin{proof}
Let  $\mu =  \delta_{0}+e^{5\pii/3}\delta_{1}+\delta_{2}+e^{2\pii/3}\delta_{3}
$.  Then 
\begin{align*} 
\mu * \tilde \mu & = \big(1+1+1+1\big)\delta_{0}
\\&\qquad
+\big(e^{4\pii/3}+e^{5\pii/3}+e^{\pii/3}+e^{2\pii/3}\big)\delta_{1}
\\&\qquad
+\big(1-1+1-1\big)\delta_{2}
\\&\qquad
+\big(e^{\pii/3}+e^{5\pii/3}+e^{4\pii/3}+e^{2\pii/3}\big)\delta_{3}\\&=4\delta_{0 }.\qedhere
\end{align*}
\end{proof}


\begin{prop}
\cite{MR627683}
\label{prop4in5}
$\{0,1,2,3\}$ is an extreme subset of $\Zp5{}$.
\end{prop}
\begin{proof}
Here is an extremal measure:
\begin{comment}\,\footnote{\ There are others. For example: $\mu = \delta_0+ e^{22\pi i/15}\delta_1 +e^{48\pi i/15}\delta_2 +e^{6\pi i/15}\delta_3.$
Then calculation shows that
\begin{align*}
\mu*\tilde\mu&= 4\delta_0
+e^{2\pi i/15}\big(-e^{\pi i/3}-e^{-\pi/3} +1\big)\delta_1
\\ &+e^{4\pi i/15}\big(1+e^{-20\pi i/15}+e^{-10\pi i/15}\big)\delta_2
\\
&+e^{\pi i/15}\big(e^{5\pi i/15}-e^{-5\pi i/15}+ e^{15\pi i/15}\big)\delta_3
\\
&+e^{3\pi i/15}\big(-e^{-10\pi i/15}-1+e^{-5\pi i/15}\big)\delta_4.
\end{align*}
The coefficient of each of $\delta_1$ to $\delta_4$  has a factor which is a sum of the third roots of three (or of their complex conjugates), which sum to zero. 
%Hence $\mu*\tilde\mu=4\delta_0$, $\|\widehat\mu\|_\infty^2=4$ and  $\mu$ is extremal.
 }%end of footnote
 \end{comment}
$\mu = \delta_0+e^{2\pi i/3}(\delta_1+\delta_2) +\delta_3$.
Then 
$
\tilde\mu = \delta_0 +e^{-2\pi i/3}(\delta_3+\delta_4)+\delta_2$, so 
\[\mu*\tilde\mu = 4\delta_0+(1+e^{2\pi i/3} +e^{-2\pi/3})(\delta_1+\delta_2+\delta_3+\delta_4)\\
= 4\delta_0. \qedhere
\]
\end{proof}

\begin{comment}
\begin{prop}\label{prop4in7}
$\{0,1,2,4\}$ is extreme in $\Zp7{}$.
\end{prop}

\end{comment}

\begin{prop}
\cite{MR627683}
\label{prop4in7}
$\{0, 1, 2, 4\}$ is extreme in \Zp{7}{}.
\end{prop}

\begin{proof}
Let $\mu = \delta_{0} - \delta_{1} - \delta_{2} - \delta_{4} 
$.  Then 
\begin{align*} 
\mu * \tilde \mu & = 4\delta_{0}   +  \big(  -1 + 1\big)\delta_{1} 
%\\&\qquad 
 +  \big(  -1 + 1\big)\delta_{2}  +  \big(  -1 + 1\big)\delta_{3} 
 \\&\qquad 
 %
  +  \big(1-1\big)\delta_{4}   +  \big(  -1 + 1\big)\delta_{5} 
   +  \big(  -1 + 1\big)\delta_{6} 
=4\delta_0.\qedhere
\end{align*}
\end{proof}


\begin{proof}[Alternative proof of \propref{prop4in7}]
Let $z =e^{2\pi i/7}$.

We compute the transform of $\mu$, using arithmetic mod 7 in
the exponents of $z$.
\begin{align*}
\widehat\mu(0) &= 1 -1-1-1=-2.\\
|\widehat\mu(1)|^2 &= -z^4 + z^3 + z^{-3} - z^{-4} + 4
\\ \notag
&= -z^4+z^3+z^4-z^3+4=4,
\\
|\widehat\mu(2)|^2 &=
-z^8 + z^6 + z^{-6} - z^{-8} + 4=4,%\text{ (further details omitted)}
\\
|\widehat\mu(3)|^2 &=
%-z^{12} + z^{9} + z^{-9} - z^{-12} + 4 \\
%\notag\quad&=
%-z^5-z^2-z^{-2}-z^{-5}+4 \\
%\notag\quad &=
-z^5-z^2-z^5-z^2+4=4,
\\
|\widehat\mu(4)|^2 &=
%-z^{16} + z^{12} + z^{-12} - z^{-16} + 4 =4,
%
-z^2+z^5+z^2-z^5+4=4,% \text{ (further details omitted)}
%\\
%|\widehat\mu(5)|^2 &=
%-z^{20} + z^{15} + z^{-15} - z^{-20} + 4=
%-z^6+z+z^{-1}-z^{-6} +4=
%-z^6+z+z^6-z+4=4, 
\text{ and}
\\
|\widehat\mu(6)|^2 &= 
%-z^{24} + z^{18} + z^{-18} - z^{-24} + 4=
%%-z^3+z^4+z^{-4}-z^{-3}+4 =
-z^3+z^4+z^3-z^4+4=4.
\qedhere
\end{align*}
\end{proof}


\bigbreak
\subsection{Sets with five elements}

\begin{prop}\label{prop5in6}
\cite{MR627683}
$\{0,1,2,3,4\}$ is extreme in \Zp{6}{}.
\end{prop}

\begin{proof} Let $\mu =  \delta_{0}  +   e^{ 3\pi i  / 2}\delta_{1}  +    -\delta_{2}  +   e^{ 3\pi i  / 2}\delta_{3}  +  \delta_{4}$.
Then
\begin{align*}
 \mu * \tilde\mu&=  5\delta_{0}   +  \big( -i - i+  i +  i \big)\delta_{1} 
  +  \big(1-1+1-1\big)\delta_{2} 
  \\& \qquad +  \big(  i - i- i+  i \big)\delta_{3} 
 % \\& \qquad
   +  \big(  -1 +1-1+ 1\big)\delta_{4} 
  \\& \qquad
   +  \big(  i +  i - i+ e^{ 3 \pi i}\big)\delta_{5}  
   \\&=5\delta_0.\qedhere
 \end{align*}
\end{proof}


\begin{prop}
\cite{MR627683}
\label{prop5in12} The subsets $E=\{0,2,4,6,8\}$, $F=\{0,2,3,4,7\}\subset \Zp{12}{}$
are not equivalent but both are extreme.
\end{prop}

\begin{remark} This does not contradict \cite{MR627683}, since one set is in a subgroup.
\end{remark}

\begin{proof} 
Non-equivalence: 
the group automorphisms of \Zp{12}{} are multiplication by 5, 7 and 11 (all mod 12). Each of them takes odd elements of \Zp{12}{} to odd elements and even elements to even elements. Translations either take evens to evens and odds to odds or evens to odds and odds to evens. Thus, no combination of group automorphisms and translations can take 
$E$, whose image will contain either only evens or only odds, onto $F$, which contains both evens and odds.

$\{0,2,4,6,8\}$ is extreme because it is a five-element subset of the the coset  $\{0,2,4,6,8,10\}$ in \Zp{12}{}.

For the second set, let $\mu =
  \delta_{0}  +   e^{ 3\pi i  / 2}\delta_{2}  +   e^{ 5\pi i  / 4}\delta_{3}  +  \delta_{4}  +   e^{\pi i / 4}\delta_{7}
.$
  Then
\begin{align*} \mu * \tilde\mu&=   5\delta_{0}   +  \big( e^{ 7 \pi i / 4}+ e^{ 3 \pi i / 4}\big)\delta_{1}
%\\&
  +  \big( -i +  i \big)\delta_{2} 
\\& \qquad
 +  \big( e^{ 5 \pi i / 4}+ e^{\pi i / 4}\big)\delta_{3} 
%\\&
  +  \big(1-1\big)\delta_{4} 
\\& \qquad +  \big( e^{ 7 \pi i / 4}+ e^{ 3 \pi i / 4}\big)\delta_{5} 
%\\&
  +  \big( e^{ 5 \pi i / 4}+ e^{\pi i / 4}\big)\delta_{7} 
\\& \qquad +  \big(1-1\big)\delta_{8} 
%\\&
  +  \big( e^{ 3 \pi i / 4}+ e^{ 7 \pi i / 4}\big)\delta_{9} 
\\& \qquad +  \big(  i - i\big)\delta_{10} 
%\\&
  +  \big( e^{\pi i / 4}+ e^{ 5 \pi i}\big)\delta_{11} 
  \\&= 5\delta_0.
\end{align*}
\end{proof}

\subsection{Sets with 6 elements}

There are no  extreme three element subsets of \Zp{10}{} nor of \Zp{14}{}, so the sets in the next two results
cannot be the sum of an extreme set with a two-element coset (all cosets having two elements in those groups).


The 6 element set $\{(0, 0),  (0, 1),  (0, 2),  (0, 3),  (0, 4),  (0, 7) \}$ is extreme in $\Zp{2}{} \times \Zp{10}{}$. \newline 
 
\begin{proof}
Let  $\mu =  \delta_{(0,0)}+e^{53\pii/30}\delta_{(0,1)}+e^{8\pii/15}\delta_{(0,2)}+e^{29\pii/30}\delta_{(0,3)}+e^{2\pii/5}\delta_{(0,4)}+e^{\pii/30}\delta_{(0,7)}
$.  Then 
\begin{align*} 
\mu * \tilde \mu & = \big(1+1+1+1+1+1\big)\delta_{(0,0)}
\\&\qquad
+\big(e^{53\pii/30}+e^{23\pii/30}+e^{13\pii/30}+e^{43\pii/30}\big)\delta_{(0,1)}
\\&\qquad
+\big(e^{8\pii/15}+e^{6\pii/5}+e^{28\pii/15}\big)\delta_{(0,2)}
\\&\qquad
+\big(e^{59\pii/30}+e^{29\pii/30}+e^{19\pii/30}+e^{49\pii/30}\big)\delta_{(0,3)}
\\&\qquad
+\big(e^{26\pii/15}+e^{2\pii/5}+e^{16\pii/15}\big)\delta_{(0,4)}
\\&\qquad
+\big(i-i\big)\delta_{(0,5)}
\\&\qquad
+\big(e^{8\pii/5}+e^{14\pii/15}+e^{4\pii/15}\big)\delta_{(0,6)}
\\&\qquad
+\big(e^{31\pii/30}+e^{41\pii/30}+e^{11\pii/30}+e^{\pii/30}\big)\delta_{(0,7)}
\\&\qquad
+\big(e^{22\pii/15}+e^{4\pii/5}+e^{2\pii/15}\big)\delta_{(0,8)}
\\&\qquad
+\big(e^{7\pii/30}+e^{37\pii/30}+e^{47\pii/30}+e^{17\pii/30}\big)\delta_{(0,9)}\\&=6\delta_{(0 , 0)}\qedhere
\end{align*}
\end{proof}


\begin{prop}\label{prop6in10} The 6-element set
$\{0,1,2,3,4,7\} $ is extreme in  \Zp{10}{}.
\end{prop}

\begin{proof}
Let $\mu =
 \delta_{0}  +   e^{ 5\pi i  / 6}\delta_{1}  +   e^{ 2\pi i  / 3}\delta_{2}  +   e^{ 5\pi i  / 6}\delta_{3}  +   e^{ 2\pi i  / 1}\delta_{4}  +   e^{\pi i / 6}\delta_{7}
.$
  Then
\begin{align*} \mu * \tilde\mu&=   6\delta_{0}   +  \big( e^{ 5 \pi i / 6}+ e^{ 11 \pi i / 6}+ e^{\pi i / 6}+ e^{ 7 \pi i / 6}\big)\delta_{1} 
  +  \big( e^{ 2 \pi i / 3}+ 1+ e^{ 4 \pi i / 3}\big)\delta_{2} 
  \\& \qquad 
  +  \big( e^{ 11 \pi i / 6}+ e^{ 5 \pi i / 6}+ e^{ 7 \pi i / 6}+ e^{\pi i / 6}\big)\delta_{3} 
  %\\& \qquad
   +  \big( e^{ 2 \pi i / 3}+ 1+ e^{ 4 \pi i / 3}\big)\delta_{4} 
  \\& \qquad
   +  \big(  i - i\big)\delta_{5} 
  %\\& \qquad
   +  \big( 1+ e^{ 2 \pi i / 3}+ e^{ 4 \pi i / 3}\big)\delta_{6} 
  \\& \qquad
   +  \big( e^{ 7 \pi i / 6}+ e^{ 5 \pi i / 6}+ e^{ 11 \pi i / 6}+ e^{\pi i / 6}\big)\delta_{7} 
 % \\& \qquad
   +  \big( e^{ 4 \pi i / 3}+ 1+ e^{ 2 \pi i / 3}\big)\delta_{8} 
  \\& \qquad +  \big( e^{ 7 \pi i / 6}+ e^{\pi i / 6}+ e^{ 11 \pi i / 6}+ e^{ 5 \pi i}\big)\delta_{9} 
  \\&=6\delta_0.
      \qedhere
\end{align*}
\end{proof}


\begin{prop}\label{prop6in14}
The set $\{ 0,1,2,3,4,7 \}$ is extreme in \Zp{14}{}.
\end{prop}

\begin{proof}
Let
$ \mu =\delta_{0}  +   e^{ 10\pi i  / 12}\delta_{1}  +   e^{ 8\pi i  / 12}\delta_{2}  +   e^{ 10\pi i  / 12}\delta_{3}  +  \delta_{4}  +   e^{ 2\pi i  / 12}\delta_{7}
.$
  Then
\begin{align*}
\mu * \tilde \mu &= 6\delta_{0}   +  \big(  i +  i - i- i\big)\delta_{1} 
%\\&
  +  \big( e^{\pi i / 1}+ 1+ e^{\pi i / 1}+ 1\big)\delta_{2} 
\\& \qquad +  \big( -i +  i +  i - i\big)\delta_{3} 
%\\&
  +  \big( 1+ e^{\pi i / 1}+ 1+ e^{\pi i / 1}\big)\delta_{4} 
\\& \qquad +  \big( -i +  i +  i - i\big)\delta_{5} 
%\\&
  +  \big( e^{\pi i / 1}+ 1+ e^{\pi i / 1}+ 1\big)\delta_{6} 
\\& \qquad +  \big(  i - i+  i +  i - i+  i - i- i\big)\delta_{7} 
\\& \qquad +  \big( e^{\pi i / 1}+ 1+ e^{\pi i / 1}+ 1\big)\delta_{8} 
%\\&
  +  \big( -i +  i +  i - i\big)\delta_{9} 
  \\&\qquad  +  \big( e^{\pi i / 1}+ 1+ e^{\pi i / 1}+ 1\big)\delta_{10} 
%\\&
  +  \big( -i - i+  i +  i \big)\delta_{11} 
\\& \qquad
  +  \big( e^{\pi i / 1}+ 1+ e^{\pi i / 1}+ 1\big)\delta_{12} 
%\\&
  +  \big( -i - i+  i + e^{\pi i / 1}\big)\delta_{13} 
  \\& = 6\delta_0.\qedhere
\end{align*}
\end{proof}

\subsection{Sets with 7 elements}\label{subsec7elts}
Seven-element exceptional extreme sets are somewhat more plentiful; we have
found one in \Zp{12}{}, two non-equivalent ones in \Zp{16}{} (and, as expected, the 7 element subset of the 8-element subgroup)
and one in \Zp{19}{}.

\begin{prop}\label{prop7in8} The 7-element set 
$\{0,1,2,3,4,5,6\}$ is extreme in \Zp8{}
\end{prop}

\begin{proof}
Let  $\mu =  \delta_{0}+e^{4\pii/3}\delta_{1}+e^{4\pii/3}\delta_{2}+e^{2\pii/3}\delta_{3}+e^{4\pii/3}\delta_{4}+e^{4\pii/3}\delta_{5}+\delta_{6}
$.  Then 
\begin{align*} 
\mu * \tilde \mu & = \big(1+1+1+1+1+1+1\big)\delta_{0}
\\&\qquad
+\big(e^{4\pii/3}+1+e^{4\pii/3}+e^{2\pii/3}+1+e^{2\pii/3}\big)\delta_{1}
\\&\qquad
+\big(1+e^{4\pii/3}+e^{4\pii/3}+1+e^{2\pii/3}+e^{2\pii/3}\big)\delta_{2}
\\&\qquad
+\big(e^{2\pii/3}+e^{4\pii/3}+e^{2\pii/3}+1+1+e^{4\pii/3}\big)\delta_{3}
\\&\qquad
+\big(e^{2\pii/3}+1+e^{4\pii/3}+e^{4\pii/3}+1+e^{2\pii/3}\big)\delta_{4}
\\&\qquad
+\big(e^{4\pii/3}+1+1+e^{2\pii/3}+e^{4\pii/3}+e^{2\pii/3}\big)\delta_{5}
\\&\qquad
+\big(e^{2\pii/3}+e^{2\pii/3}+1+e^{4\pii/3}+e^{4\pii/3}+1\big)\delta_{6}
\\&\qquad
+\big(e^{2\pii/3}+1+e^{2\pii/3}+e^{4\pii/3}+1+e^{4\pii/3}\big)\delta_{7}\\&=7\delta_{(0 , 0}\qedhere
\end{align*}
\end{proof}
\begin{comment}
\begin{proof}
Let 
$\nu=\delta_0+
e^{11\pi i/6}\delta_1+
e^{14\pi i/6} \delta_2
+e^{13\pi i/6}\delta_3
+e^{8\pi i/6}\delta_4
+e^{11\pi i/6}\delta_5
+e^{6\pi i/6}\delta_6.
$
Then
\begin{align*}
\nu*\widetilde\nu &= 7\delta_0 +
\big( e^{ 11 \pi i / 6}+  i + e^{ 11 \pi i / 6}+ e^{ 7 \pi i / 6}+  i + e^{ 7 \pi i / 6}\big)\delta_{1} 
\\& \qquad +  \big( e^{pi i}+ e^{\pi i / 3}+ e^{\pi i / 3}+ e^{pi i}+ e^{ 5 \pi i / 3}+ e^{ 5 \pi i / 3}\big)\delta_{2} 
\\& \qquad +  \big( e^{\pi i / 6}+ e^{ 5 \pi i / 6}+ e^{\pi i / 6}- i- i+ e^{ 5 \pi i / 6}\big)\delta_{3}
\\& \qquad +  \big( e^{ 2 \pi i / 3}+ 1+ e^{ 4 \pi i / 3}+ e^{ 4 \pi i / 3}+ 1+ e^{ 2 \pi i / 3}\big)\delta_{4} 
\\& \qquad +  \big( e^{ 11 \pi i / 6}+  i +  i + e^{ 7 \pi i / 6}+ e^{ 11 \pi i / 6}+ e^{ 7 \pi i / 6}\big)\delta_{5}
\\& \qquad +  \big( e^{ 5 \pi i / 3}+ e^{ 5 \pi i / 3}+ e^{pi i}+ e^{\pi i / 3}+ e^{\pi i / 3}- 1 \big)\delta_{6} 
\\& \qquad +  \big( e^{\pi i / 6}- i+ e^{\pi i / 6}+ e^{ 5 \pi i / 6}- i+ e^{ 5 \pi i}\big)\delta_{7}
\\&=7\delta_0. \qedhere
\end{align*}
\end{comment}
\begin{comment}\footnote{\ An extreme measure on ${\0,\dots,5\}\subset\Zp8{}$
is $\mu = \delta_0 +e^{4\pi i/3}\delta_1+e^{4\pi i/3}\delta_2
+e^{2\pi i/3}\delta_3 + e^{4\pi i/3}\delta_4+e^{4\pi i/3}\delta_5+\delta_6$.
Indeed,
\begin{align*}
\mu * \tilde \mu  = 7\delta_0 &+\prnb{e^{4\pi i/3} + e^{-4\pi/3}+1+e^{2\pi i/3}+e^{-2\pi i/3} +1}\delta_1 \\
&+\prnb{ 1+e^{-4\pi i/3}+e^{2\pi i/3} +1 + e^{-2\pi i/3}+e^{4\pi i/3} }\delta_2 \\
&+\prnb{e^{2\pi i/3} +e^{-4\pi i/3} + e^{-2\pi i/3} +1+e^{4\pi i/3}  +1}\delta_3 \\
&+\prnb{ e^{4\pi i/3} + e^{4\pi i/3} +1 + e^{-4\pi i/3} + e^{-4\pi i/3} +1 }\delta_4 \\
&\cdots\\
&=7\delta_0.
\end{align*}
}
\end{proof}
\end{comment}


\begin{rem}
Extreme measures are not unique, even when they are required to have mass 1 at the identity.
 Another extreme measure on $\{0,\dots,6\}\subset \Zp{8}{}$
is \[
\nu=\delta_0+
e^{11\pi i/6}\delta_1+
e^{14\pi i/6} \delta_2
+e^{13\pi i/6}\delta_3
+e^{8\pi i/6}\delta_4
+e^{11\pi i/6}\delta_5
+e^{6\pi i/6}\delta_6.
\]
We omit the proof that $\nu*\widetilde{\nu} =7\delta_0.$
\begin{comment}
Indeed,
\begin{align*}
\nu*\widetilde\nu &= 7\delta_0 +
\big( e^{ 11 \pi i / 6}+  i + e^{ 11 \pi i / 6}+ e^{ 7 \pi i / 6}+  i + e^{ 7 \pi i / 6}\big)\delta_{1} 
\\& \qquad +  \big( e^{pi i}+ e^{\pi i / 3}+ e^{\pi i / 3}+ e^{pi i}+ e^{ 5 \pi i / 3}+ e^{ 5 \pi i / 3}\big)\delta_{2} 
\\& \qquad +  \big( e^{\pi i / 6}+ e^{ 5 \pi i / 6}+ e^{\pi i / 6}- i- i+ e^{ 5 \pi i / 6}\big)\delta_{3}
\\& \qquad +  \big( e^{ 2 \pi i / 3}+ 1+ e^{ 4 \pi i / 3}+ e^{ 4 \pi i / 3}+ 1+ e^{ 2 \pi i / 3}\big)\delta_{4} 
\\& \qquad +  \big( e^{ 11 \pi i / 6}+  i +  i + e^{ 7 \pi i / 6}+ e^{ 11 \pi i / 6}+ e^{ 7 \pi i / 6}\big)\delta_{5}
\\& \qquad +  \big( e^{ 5 \pi i / 3}+ e^{ 5 \pi i / 3}+ e^{pi i}+ e^{\pi i / 3}+ e^{\pi i / 3}- 1 \big)\delta_{6} 
\\& \qquad +  \big( e^{\pi i / 6}- i+ e^{\pi i / 6}+ e^{ 5 \pi i / 6}- i+ e^{ 5 \pi i}\big)\delta_{7}
\\&=7\delta_0. \qedhere
\end{align*}
\end{comment}
\end{rem}

\begin{prop}\label{prop7in12} The 7-element set
$\{0,1,2,5,6,8,9\}$ is extreme in \Zp{12}{}.
\end{prop}

\begin{proof} Let $\mu =\delta_{0}  +   e^{ 7\pi i  / 12}\delta_{1}  +  \delta_{3}  +   e^{ 5\pi i  / 6}\delta_{4}  +    i \delta_{6}  +   e^{ 7\pi i  / 12}\delta_{7}  +    i \delta_{9}  +   e^{ 5\pi i  / 6}\delta_{10}
$.
Then
\begin{align*} \mu*\tilde \mu  &=7\delta_{0}   +  \big( e^{ 7 \pi i / 12}+ e^{ 5 \pi i / 6}+ e^{\pi i / 12}+ e^{\pi i / 3}\big)\delta_{1} \\& \qquad +  \big( e^{ 7 \pi i / 6}+ e^{ 17 \pi i / 12}+ e^{ 5 \pi i / 3}+ e^{ 23 \pi i / 12}\big)\delta_{2} \\& \qquad +  \big( -i + e^{ 7 \pi i / 4}+ 1+ e^{\pi i / 4}+  i + e^{ 7 \pi i / 4}+ 1+ e^{\pi i / 4}\big)\delta_{3} \\& \qquad +  \big( e^{\pi i / 12}+ e^{ 5 \pi i / 6}+ e^{ 7 \pi i / 12}+ e^{\pi i / 3}\big)\delta_{4} \\& \qquad +  \big( e^{ 17 \pi i / 12}+ e^{ 7 \pi i / 6}+ e^{ 23 \pi i / 12}+ e^{ 5 \pi i / 3}\big)\delta_{5} \\& \qquad +  \big( -i + 1- i+ 1+  i + 1+  i + 1\big)\delta_{6} \\& \qquad +  \big( e^{\pi i / 12}+ e^{\pi i / 3}+ e^{ 7 \pi i / 12}+ e^{ 5 \pi i / 6}\big)\delta_{7} \\& \qquad +  \big( e^{ 7 \pi i / 6}+ e^{ 17 \pi i / 12}+ e^{ 5 \pi i / 3}+ e^{ 23 \pi i / 12}\big)\delta_{8} \\& \qquad +  \big( 1+ e^{ 7 \pi i / 4}- i+ e^{\pi i / 4}+ 1+ e^{ 7 \pi i / 4}+  i + e^{\pi i / 4}\big)\delta_{9} \\& \qquad +  \big( e^{ 7 \pi i / 12}+ e^{\pi i / 3}+ e^{\pi i / 12}+ e^{ 5 \pi i / 6}\big)\delta_{10} \\& \qquad +  \big( e^{ 17 \pi i / 12}+ e^{ 7 \pi i / 6}+ e^{ 23 \pi i / 12}+ e^{ 5 \pi i}\big)\delta_{11} 
\\&=7\delta_0. 
\end{align*}
Each sum in the parenthese above is zero.
\begin{comment}\footnote{ Another extreme measure is: $ \mu =\delta_{0}  +   e^{ 17\pi i  / 12}\delta_{1}  +    i \delta_{2}  +   e^{ 13\pi i  / 12}\delta_{5}  +   e^{ 3\pi i  / 2}\delta_{6}  +    -\delta_{8}  +   e^{ 13\pi i  / 12}\delta_{9}
.$
Then
\begin{align*} 
\mu*\tilde \mu =7\delta_{0}&  +  \big( e^{ 17 \pi i / 12}+ e^{ 13 \pi i / 12}+ e^{ 5 \pi i / 12}+ e^{\pi i / 12}\big)\delta_{1}   +  \big(  i - i\big)\delta_{2} 
\\&  +  \big( e^{ 11 \pi i / 12}+ e^{ 7 \pi i / 12}+ e^{ 23 \pi i / 12}+ e^{ 19 \pi i / 12}\big)\delta_{3}   +  \big( e^{pi i}+ e^{\pi i / 3}+ e^{ 5 \pi i / 3}+ e^{pi i}+ 1\big)\delta_{4} 
\\&  +  \big( e^{ 5 \pi i / 12}+ e^{ 17 \pi i / 12}+ e^{ 13 \pi i / 12}+ e^{\pi i / 12}\big)\delta_{5}   +  \big(  i - i- i+  i \big)\delta_{6} 
\\&  +  \big( e^{ 11 \pi i / 12}+ e^{ 23 \pi i / 12}+ e^{ 19 \pi i / 12}+ e^{ 7 \pi i / 12}\big)\delta_{7}   +  \big( e^{\pi i / 3}+ e^{pi i}+ 1+ e^{pi i}+ e^{ 5 \pi i / 3}\big)\delta_{8} 
\\&  +  \big( e^{ 17 \pi i / 12}+ e^{\pi i / 12}+ e^{ 5 \pi i / 12}+ e^{ 13 \pi i / 12}\big)\delta_{9}   +  \big( -i +  i \big)\delta_{10} 
\\&  +  \big( e^{ 7 \pi i / 12}+ e^{ 11 \pi i / 12}+ e^{ 19 \pi i / 12}+ e^{ 23 \pi i}\big)\delta_{11}\\&
= 7\delta_0.
\end{align*}
} %end footnote
\end{comment}

\qedhere
\end{proof}



\begin{prop}\label{prop7in16} Each of 7-element sets
\begin{enumerate}
\item
 $\{0,2,4,6,8,10,12\}$ and 
 \item $\{0,1,2,4,5,7,11\}$  
 \end{enumerate}
 is extreme in \Zp{16}{}.
They are not equivalent.
\end{prop}

\begin{proof}
The sets are  not equivalent because the automorphisms 
(multiplication by odd integers) preserve the parity  of
 elements and translation switches or leaves fixed the 
 parity. Since the first set has both odd and even elements
 so will every set equivalent to it. (We will use this argument
 again in \propref{prop8in16}.)
 
 As for extremality, first note that $\{0,2,4,6,8,10,12\}$ is extreme because it is a 7 element subset of an 8 element coset
and so extreme by \propref{prop7in8}. For the other set, set
 \[\mu = 
  \delta_{0}     -\delta_{1}  +   e^{ 4\pi i  / 3}\delta_{2}  +  \delta_{4}  +   e^{ 5\pi i  / 3}\delta_{5}  +   e^{ 5\pi i  / 3}\delta_{7}  +    -\delta_{11}.
\]
Then
\begin{align*}  \mu*\tilde \mu = 7\delta_{0} &  + \big(  -1 + e^{\pi i / 3}+ e^{ 5 \pi i / 3}\big)\delta_{1}   +  \big( e^{ 4 \pi i / 3}+ e^{ 2 \pi i / 3}+ 1\big)\delta_{2} \\&  +  \big(  -1 + e^{\pi i / 3}+ e^{ 5 \pi i / 3}\big)\delta_{3}   +  \big( 1+ e^{ 2 \pi i / 3}+ e^{ 4 \pi i / 3}\big)\delta_{4} \\&  +  \big(  -1 + e^{ 5 \pi i / 3}+ e^{\pi i / 3}\big)\delta_{5}   +  \big( 1+ e^{ 2 \pi i / 3}+ e^{ 4 \pi i / 3}\big)\delta_{6} \\&  +  \big( e^{\pi i / 3}+ e^{ 5 \pi i / 3}- 1 \big)\delta_{7}   +  \big( e^{\pi i / 3}- 1 + e^{ 5 \pi i / 3}\big)\delta_{9} \\&  +  \big( e^{ 4 \pi i / 3}+ e^{ 2 \pi i / 3}+ 1\big)\delta_{10}   +  \big( e^{\pi i / 3}+ e^{ 5 \pi i / 3}- 1 \big)\delta_{11} \\&  +  \big( 1+ e^{ 4 \pi i / 3}+ e^{ 2 \pi i / 3}\big)\delta_{12}   +  \big(  -1 + e^{ 5 \pi i / 3}+ e^{\pi i / 3}\big)\delta_{13} \\&  +  \big( e^{ 2 \pi i / 3}+ e^{ 4 \pi i / 3}+ 1\big)\delta_{14}   +  \big(  -1 + e^{ 5 \pi i / 3}- 1 \big)\delta_{15} \\&=7\delta_0.
\qedhere
\end{align*}
\end{proof}

\begin{prop}\label{prop7in19} The 7-element set
$\{0,1,2,5,12,13,15\}$ is extreme in \Zp{19}{}.
\end{prop}

\begin{proof} 
Let $\mu =\delta_{0}  +   e^{ 4\pi i  / 3}\delta_{1}  +  \delta_{2}  +  \delta_{5}  +   e^{ 4\pi i  / 3}\delta_{12}  +   e^{ 4\pi i  / 3}\delta_{13}  +   e^{\pi i / 3}\delta_{15}.$
Then
\begin{align*} \mu*\tilde \mu &=
7\delta_{0}   +  \big( e^{ 4 \pi i / 3}+ e^{ 2 \pi i / 3}+ 1\big)\delta_{1} 
  +  \big( 1+ e^{\pi i }\big)\delta_{2} 
  \\& \qquad +  \big( 1+ e^{\pi i }\big)\delta_{3} 
%  \\& \qquad
   +  \big( e^{ 5 \pi i / 3}+ e^{ 2 \pi i / 3}\big)\delta_{4} \\& \qquad +  \big( e^{\pi i }+ 1\big)\delta_{5} 
  %\\& \qquad
   +  \big( e^{ 2 \pi i / 3}+ e^{ 5 \pi i / 3}\big)\delta_{6} \\& \qquad +  \big( e^{ 2 \pi i / 3}+ 1+ e^{ 4 \pi i / 3}\big)\delta_{7} 
  % \\& \qquad
    +  \big( 1+ e^{ 2 \pi i / 3}+ e^{ 4 \pi i / 3}\big)\delta_{8} \\& \qquad +  \big( e^{ 2 \pi i / 3}+ e^{ 5 \pi i / 3}\big)\delta_{9}
 %    \\& \qquad
      +  \big( e^{ 4 \pi i / 3}+ e^{\pi i / 3}\big)\delta_{10} \\& \qquad +  \big( e^{ 2 \pi i / 3}+ 1+ e^{ 4 \pi i / 3}\big)\delta_{11}
  %     \\& \qquad
 +  \big( e^{ 2 \pi i / 3}+ e^{ 4 \pi i / 3}+ 1\big)\delta_{12} \\& \qquad +  \big( e^{ 4 \pi i / 3}+ e^{\pi i / 3}\big)\delta_{13} 
 %\\&
   +  \big( 1+ e^{\pi i }\big)\delta_{14} \\& \qquad +  \big( e^{ 4 \pi i / 3}+ e^{\pi i / 3}\big)\delta_{15} 
     +  \big( 1+ e^{\pi i }\big)\delta_{16} \\& \qquad +  \big( 1+ e^{\pi i }\big)\delta_{17} 
       +  \big( e^{ 2 \pi i / 3}+ e^{ 4 \pi i / 3}+ 1\big)\delta_{18} 
       \\&= 7\delta_0.  \qedhere
\end{align*}
\end{proof}

We have found no other 7-element extreme sets in cyclic groups of order below 20.
\subsection{Sets with 8 elements}
\begin{prop}\label{prop8inZ9}
The 8 element set $\{0,  1,  2,  3,  4,  5,  6,  7 \}$ is extreme in $\Zp{9}{}$. \newline \end{prop} 


 
\begin{proof}[First proof]
Let  $\mu =  \delta_{0}+e^{10\pii/7}\delta_{1}+e^{4\pii/7}\delta_{2}+e^{2\pii/7}\delta_{3}+e^{2\pii/7}\delta_{4}+e^{4\pii/7}\delta_{5}+e^{10\pii/7}\delta_{6}+\delta_{7}
$.  Then 
\begin{align*} 
\mu * \tilde \mu & = \big(1+1+1+1+1+1+1+1\big)\delta_{0}
\\&\qquad
+\big(e^{10\pii/7}+e^{8\pii/7}+e^{12\pii/7}+1+e^{2\pii/7}+e^{6\pii/7}+e^{4\pii/7}\big)\delta_{1}
\\&\qquad
+\big(1+e^{4\pii/7}+e^{6\pii/7}+e^{12\pii/7}+e^{2\pii/7}+e^{8\pii/7}+e^{10\pii/7}\big)\delta_{2}
\\&\qquad
+\big(e^{4\pii/7}+e^{10\pii/7}+e^{2\pii/7}+e^{6\pii/7}+1+e^{8\pii/7}+e^{12\pii/7}\big)\delta_{3}
\\&\qquad
+\big(e^{10\pii/7}+1+e^{4\pii/7}+e^{2\pii/7}+e^{8\pii/7}+e^{6\pii/7}+e^{12\pii/7}\big)\delta_{4}
\\&\qquad
+\big(e^{12\pii/7}+e^{6\pii/7}+e^{8\pii/7}+e^{2\pii/7}+e^{4\pii/7}+1+e^{10\pii/7}\big)\delta_{5}
\\&\qquad
+\big(e^{12\pii/7}+e^{8\pii/7}+1+e^{6\pii/7}+e^{2\pii/7}+e^{10\pii/7}+e^{4\pii/7}\big)\delta_{6}
\\&\qquad
+\big(e^{10\pii/7}+e^{8\pii/7}+e^{2\pii/7}+e^{12\pii/7}+e^{6\pii/7}+e^{4\pii/7}+1\big)\delta_{7}
\\&\qquad
+\big(e^{4\pii/7}+e^{6\pii/7}+e^{2\pii/7}+1+e^{12\pii/7}+e^{8\pii/7}+e^{10\pii/7}\big)\delta_{8}\\&=8\delta_{0 }\qedhere
\end{align*}
\end{proof}

 

\begin{proof}[Second proof] Let \begin{align*}\mu =
 \delta_{0}  + &  e^{ 20\pi i  / 21}\delta_{1}  +   e^{ 64\pi i  / 21}\delta_{2}  +   e^{ 20\pi i  / 7}\delta_{3}  +   e^{ 32\pi i  / 21}\delta_{4}  \\ &\qquad +   e^{ 64\pi i  / 21}\delta_{5}  +   e^{ 16\pi i  / 7}\delta_{6}  +   e^{ 8\pi i  / 3}\delta_{7}
.\end{align*}
  Then
\begin{align*} \mu * \tilde\mu&=   8\delta_{0}  +  \big( e^{ 20 \pi i / 21}+ e^{ 2 \pi i / 21}+ e^{ 38 \pi i / 21}+ e^{ 2 \pi i / 3}+ e^{ 32 \pi i / 21}
\\& \qquad \qquad+ e^{ 26 \pi i / 21}+ e^{ 8 \pi i / 21}\big)\delta_{1} 
\\& \qquad
 +  \big( e^{ 4 \pi i / 3}+ e^{ 22 \pi i / 21}+ e^{ 40 \pi i / 21}+ e^{ 10 \pi i / 21}
\\& \qquad \qquad+ e^{ 4 \pi i / 21}+ e^{ 16 \pi i / 21}+ e^{ 34 \pi i / 21}\big)\delta_{2} 
\\& \qquad
 +  \big( e^{ 12 \pi i / 7}+ e^{ 2 \pi i / 7}+ e^{ 6 \pi i / 7}+ e^{ 4 \pi i / 7}
\\& \qquad \qquad+ 1+ e^{ 10 \pi i / 7}+ e^{ 8 \pi i / 7}\big)\delta_{3} 
\\& \qquad
 +  \big( e^{ 20 \pi i / 21}+ e^{ 2 \pi i / 3}+ e^{ 8 \pi i / 21}+ e^{ 32 \pi i / 21}
\\& \qquad \qquad+ e^{ 2 \pi i / 21}+ e^{ 26 \pi i / 21}+ e^{ 38 \pi i / 21}\big)\delta_{4} 
\\& \qquad
 +  \big( e^{ 10 \pi i / 21}+ e^{ 40 \pi i / 21}+ e^{ 16 \pi i / 21}+ e^{ 4 \pi i / 21}
\\& \qquad \qquad+ e^{ 22 \pi i / 21}+ e^{ 4 \pi i / 3}+ e^{ 34 \pi i / 21}\big)\delta_{5} 
\\& \qquad
 +  \big( e^{ 8 \pi i / 7}+ e^{ 10 \pi i / 7}+ 1+ e^{ 4 \pi i / 7}+ e^{ 6 \pi i / 7}+ e^{ 2 \pi i / 7}+ e^{ 12 \pi i / 7}\big)\delta_{6} 
\\& \qquad
 +  \big( e^{ 20 \pi i / 21}+ e^{ 2 \pi i / 21}+ e^{ 32 \pi i / 21}+ e^{ 38 \pi i / 21}
\\& \qquad \qquad+ e^{ 26 \pi i / 21}+ e^{ 8 \pi i / 21}+ e^{ 2 \pi i / 3}\big)\delta_{7} 
\\& \qquad +  \big( e^{ 22 \pi i / 21}+ e^{ 40 \pi i / 21}+ e^{ 4 \pi i / 21}
\\& \qquad \qquad+ e^{ 4 \pi i / 3}+ e^{ 10 \pi i / 21}+ e^{ 16 \pi i / 21}+ e^{ 34 \pi i}\big)\delta_{8} 
\\&= 9\delta_0.
      \qedhere
\end{align*}
\end{proof}




%16	8	 {0 1 4 5 8 9 12 13 }	2.828427125	2.828427125	187	12	  0 11 9 11 0 5 9 5

%16	8	 {0 2 4 6 8 10 12 14 }	2.828427125	2.828427125	26	12	  0 9 6 6 6 9 0 6
%Autos are mult by 1,3,5,7,9, 11,13,15
%and $\{0, 1, 4, 5, 8, 9, 12, 13\}$ are


\begin{prop}\label{prop8in10} The 8-element set
$\{0,1,2,3,5,6,7,8\}$ is extreme in \Zp{10}{}. %They are not equivalent.
\end{prop}


\begin{proof}[Computational proof of \propref{prop8in10}] 
Let $\mu=\delta_{0}  +   e^{ 7\pi i  / 6}\delta_{1}  +   e^{ 2\pi i  / 3}\delta_{2}  +    i \delta_{3}  +    i \delta_{5}  +   e^{ 2\pi i  / 3}\delta_{6}  +   e^{ 7\pi i  / 6}\delta_{7}  +  \delta_{8} $.
Then
\begin{align*}  
\mu*\tilde\mu &=   8\delta_{0}  +  \big( e^{ 7 \pi i / 6}- i+ e^{ 11 \pi i / 6}+ e^{\pi i / 6}+  i + e^{ 5 \pi i / 6}\big)\delta_{1} 
\\& \qquad +  \big( 1+ e^{ 2 \pi i / 3}+ e^{ 4 \pi i / 3}+ 1+ e^{ 2 \pi i / 3}+ e^{ 4 \pi i / 3}\big)\delta_{2} 
\\& \qquad +  \big( e^{ 5 \pi i / 6}+ e^{ 7 \pi i / 6}+  i + e^{ 11 \pi i / 6}+ e^{\pi i / 6}- i\big)\delta_{3} 
\\& \qquad +  \big( e^{ 4 \pi i / 3}+ 1+ e^{ 2 \pi i / 3}+ e^{ 4 \pi i / 3}+ 1+ e^{ 2 \pi i / 3}\big)\delta_{4} 
\\& \qquad +  \big( -i +  i - i+  i +  i - i+  i - i\big)\delta_{5} 
\\& \qquad +  \big( e^{ 2 \pi i / 3}+ 1+ e^{ 4 \pi i / 3}+ e^{ 2 \pi i / 3}+ 1+ e^{ 4 \pi i / 3}\big)\delta_{6} 
\\& \qquad +  \big( -i + e^{\pi i / 6}+ e^{ 11 \pi i / 6}+  i + e^{ 7 \pi i / 6}+ e^{ 5 \pi i / 6}\big)\delta_{7} 
\\& \qquad +  \big( e^{ 4 \pi i / 3}+ e^{ 2 \pi i / 3}+ 1+ e^{ 4 \pi i / 3}+ e^{ 2 \pi i / 3}+ 1\big)\delta_{8} 
\\& \qquad +  \big( e^{ 5 \pi i / 6}+  i + e^{\pi i / 6}+ e^{ 11 \pi i / 6}- i+ e^{ 7 \pi i }\big)\delta_{9}
\\&=9\delta_0.\qedhere
\end{align*}
\end{proof}


\begin{proof}[Direct proof of \propref{prop8in10}] The set is the sum of a 4-element set in the 5-element
subgroup $\{0,2,4,6,8\}$ of \Zp{10}{} and the subgroup $\{0,5\}$:
$\{0,1,2,3,5,6,7,8\} = \{0,2,6,8\}+\{0,5\}=\{0,2,6,8,5,7,11,13\} = \{0,2,6,8,5,7,1,3\}$ $\mod 10$.
\end{proof}

Turning to \Zp{12}{}, we observe that an 8-element set there could be a sum of a 4-element set and a 2-element set.
However, we cannot get an extreme set that way, though there is an 8-element extreme set in \Zp{12}{}, as the
following Lemma and Proposition show.

\begin{lem}\label{lemsumsinZtwlv}
\Zp{12}{} has no 8-element set that is a sum of a four-element extreme set and a coset.
\end{lem}

\begin{proof}
Suppose $E= A+B$ has 8 elements, where $A$ has 4 elements and $B$ has two. We may assume $0\in A$ and $B=\{0,6\}$.
Clearly $A$ cannot contain 6. \refitem{remitOneRepGen} and calculation (both by hand and machine) show that if a 4-element set in \Zp{12}{} 
 lacks 6, it is not 
extreme.\footnote{\,
The only 4-element extreme sets of \Zp{12}{} (up to equivalence) 
are $\{0,1,6,7\},$ $\{0, 2,6,8\}$ and $\{0,3,6,9\}$, all containing 6. Several other four-element sets pass the test of
\refitem{remitOneRepGen} but are not extreme.}% end footnote
%%$\{0,1,2,3\}+\{0,6\}= \{0,1,2,3,6,7,8,9\}$
\end{proof}


\begin{prop}\label{prop8in12} The 8-element set
$\{0,1,3,4,6,7,9,10\}$ % $ \{0,1,3,4\} +\{0,6\}$  
 is extreme in \Zp{12}{} and is not the sum of two extreme sets in spite of being the sum $ \{0,1,3,4\} +\{0,6\}$
 and also the sum $\{0,3,6,9\}+\{0,1\}$. 
\end{prop}
\begin{proof}Let $\mu =  \delta_{0}  +   e^{ 7\pi i  / 6}\delta_{1}  +  \delta_{3}  +   e^{ 5\pi i  / 3}\delta_{4}  +   e^{\pi i }\delta_{6}  +   e^{ 7\pi i  / 6}\delta_{7}  +   e^{\pi i }\delta_{9}  +   e^{ 5\pi i  / 3}\delta_{10}$.

Indeed, each of the sums in parentheses below is a net zero:
\begin{align*} \mu*\tilde \mu = 8\delta_{0} &   +  \big( e^{ 7 \pi i / 6}+ e^{ 5 \pi i / 3}+ e^{\pi i / 6}+ e^{ 2 \pi i / 3}\big)\delta_{1} 
\\&
 +  \big( e^{\pi i / 3}+ e^{ 5 \pi i / 6}
+ e^{ 4 \pi i / 3}+ e^{ 11 \pi i / 6}\big)\delta_{2} 
\\
&  +  \big( e^{\pi i }- i+ 1+  i + e^{\pi i }- i+ 1+  i \big)\delta_{3} 
\\
&  +  \big( e^{\pi i / 6}+ e^{ 5 \pi i / 3}+ e^{ 7 \pi i / 6}+ e^{ 2 \pi i / 3}\big)\delta_{4} 
\\&
  +  \big( e^{ 5 \pi i / 6}+ e^{\pi i / 3}+ e^{ 11 \pi i / 6}+ e^{ 4 \pi i / 3}\big)\delta_{5} 
\\
&  +  \big( e^{\pi i }+ 1+ e^{\pi i }+ 1+ e^{\pi i }+ 1+ e^{\pi i }+ 1\big)\delta_{6} 
\\&
  +  \big( e^{\pi i / 6}+ e^{ 2 \pi i / 3}+ e^{ 7 \pi i / 6}+ e^{ 5 \pi i / 3}\big)\delta_{7} 
\\
&  +  \big( e^{\pi i / 3}+ e^{ 5 \pi i / 6}+ e^{ 4 \pi i / 3}+ e^{ 11 \pi i / 6}\big)\delta_{8} 
\\
&  +  \big( 1- i+ e^{\pi i }+  i + 1- i+ e^{\pi i }+  i \big)\delta_{9} 
\\&
  +  \big( e^{ 7 \pi i / 6}+ e^{ 2 \pi i / 3}+ e^{\pi i / 6}+ e^{ 5 \pi i / 3}\big)\delta_{10} 
\\&
  +  \big( e^{ 5 \pi i / 6}+ e^{\pi i / 3}+ e^{ 11 \pi i / 6}+ e^{ 4 \pi i }\big)\delta_{11}
  \\& = 8\delta_0. 
\end{align*}
Now apply \lemref{lemsumsinZtwlv} and the fact that $\{0,1\}$ is not extreme in \Zp{12}{}\footnote{The difference 
$\{0,1\}-\{0,1\}$ produces the terms $0-0$, 1-1, 0-1, 1-0 and so $\{0,1\} \subset \Zp{12}{} $ fails the test of \refitem{remitOneRepGen}.}
\end{proof}



%16	8	 {0 1 4 5 8 9 12 13 }	2.828427125	2.828427125	187	12	  0 11 9 11 0 5 9 5

%16	8	 {0 2 4 6 8 10 12 14 }	2.828427125	2.828427125	26	12	  0 9 6 6 6 9 0 6
%Autos are mult by 1,3,5,7,9, 11,13,15
\begin{prop}\label{prop8in16} Each of the 8-element sets
\begin{enumerate}
\item
$\{0, 1, 4, 5, 8, 9, 12, 13\}$ 
\item $\{0,2,4,6,8,10,12,14\}$ 
\end{enumerate}
is  extreme in \Zp{16}{}. They are not equivalent.
\end{prop}

\begin{proof}
The sets are  not equivalent because the autormorphisms 
(multiplication by odd integers) preserve the parity  of
 elements and translation switches or leaves fixed the 
 parity. Since the first set has both odd and even elements, 
 so will every set equivalent to it. 
 
 The second set is extreme because it is a coset.
 
For the extremality of the first set, let $\mu = \delta_{0} + e^{ 22 \pi i / 12}\delta_{1} + e^{ 18 \pi i / 12}\delta_{4} + e^{ 22 \pi i / 12}\delta_{5} +\delta_{8} + e^{ 10 \pi i / 12}\delta_{9} + e^{ 18 \pi i / 12}\delta_{12} + e^{ 10 \pi i / 12}\delta_{13} 
$.  Then 
\begin{align*} 
\mu * \tilde \mu & = 8\delta_{0}   +  \big( e^{ 11 \pi i / 6}+ e^{\pi i / 3}+ e^{ 5 \pi i / 6}+ e^{ 4 \pi i / 3}\big)\delta_{1}
\\&\qquad  
+  \big( e^{ 7 \pi i / 6}+ e^{ 5 \pi i / 3}+ e^{\pi i / 6}+ e^{ 2 \pi i / 3}\big)\delta_{3} \\&\qquad  +  \big( e^{\pi i / 2}+ e^{\pi i}+ e^{ 3 \pi i / 2}+ 1+ e^{\pi i / 2}+ e^{\pi i}+ e^{ 3 \pi i / 2}+ 1\big)\delta_{4} \\&\qquad  +  \big( e^{\pi i / 3}+ e^{ 11 \pi i / 6}+ e^{ 4 \pi i / 3}+ e^{ 5 \pi i / 6}\big)\delta_{5}
 \\&\qquad  
 +  \big( e^{ 7 \pi i / 6}+ e^{ 2 \pi i / 3}+ e^{\pi i / 6}+ e^{ 5 \pi i / 3}\big)\delta_{7} \\&\qquad  +  \big( 1+ e^{\pi i}+ 1+ e^{\pi i}+ 1+ e^{\pi i}+ 1+ e^{\pi i}\big)\delta_{8} \\&\qquad  +  \big( e^{ 11 \pi i / 6}+ e^{\pi i / 3}+ e^{ 5 \pi i / 6}+ e^{ 4 \pi i / 3}\big)\delta_{9} \\&\qquad  +  \big( e^{\pi i / 6}+ e^{ 2 \pi i / 3}+ e^{ 7 \pi i / 6}+ e^{ 5 \pi i / 3}\big)\delta_{11} \\&\qquad  +  \big( e^{\pi i / 2}+ 1+ e^{ 3 \pi i / 2}+ e^{\pi i}+ e^{\pi i / 2}+ 1+ e^{ 3 \pi i / 2}+ e^{\pi i}\big)\delta_{12} \\&\qquad  +  \big( e^{\pi i / 3}+ e^{ 11 \pi i / 6}+ e^{ 4 \pi i / 3}+ e^{ 5 \pi i / 6}\big)\delta_{13} \\&\qquad  +  \big( e^{\pi i / 6}+ e^{ 5 \pi i / 3}+ e^{ 7 \pi i / 6}+ e^{ 2 \pi i}\big)\delta_{15} 
\\&=8\delta_0.\qedhere
\end{align*}
\end{proof}



\break


\subsection{Sets with 9 elements}
 
  
\begin{prop}\label{prop9in10}
$E=\{0,1,2,3,4,5,6,7,8\}$ is extreme in \Zp{10}{}.
\end{prop}

\begin{proof} Let $\mu =
 \delta_{0}  +   e^{ 3\pi i  / 2}\delta_{1}  +  \delta_{2}  +    i \delta_{3}  +    -\delta_{4}  +    i \delta_{5}  +  \delta_{6}  +   e^{ 3\pi i  / 2}\delta_{7}  +  \delta_{8}
.$
  Then
\begin{align*} \mu * \tilde\mu&=   9\delta_{0}     +  \big( -i +  i +  i +  i - i- i- i+  i \big)\delta_{1} \\& \qquad +  \big( 1+1-1- 1 +1-1- 1 + 1\big)\delta_{2} \\& \qquad +  \big(  i - i+  i - i+  i - i+  i - i\big)\delta_{3} \\& \qquad +  \big( 1+ 1+1-1- 1 +1-1- 1 \big)\delta_{4} \\& \qquad +  \big( -i - i+  i +  i +  i +  i - i- i\big)\delta_{5} \\& \qquad +  \big(  -1 - 1 +1-1- 1 + 1+ 1+ 1\big)\delta_{6} \\& \qquad +  \big( -i +  i - i+  i - i+  i - i+  i \big)\delta_{7} \\& \qquad +  \big(1-1- 1 +1-1- 1 + 1+ 1\big)\delta_{8} \\& \qquad +  \big(  i - i- i- i+  i +  i +  i + e^{ 3 \pi i}\big)\delta_{9}
\\&= 9\delta_0.
\end{align*}

\end{proof}



\begin{prop}\label{prop9in12}
$\{0,1,2,3,4,5,6,7,8\}$ and $\{0,1,2,4,5,6,8,9,10\}$ are extreme in \Zp{12}{}. They are  non-equivalent. 
\end{prop}

\begin{proof}
There are two ways to prove non-equivalence. First, by  a tedious calculation (which we delegated to a computer). The second is to show that one of the sets is the sum
of a coset with a three-element set and the other is not, which we do in the next paragraph.


The only 3-element subgroup in \Zp{12}{} is $\{0,4,8\}$. If $\{0,4,8\}+\{a\} 
\subset\{0,1,2,3,4,5,6,7,8\}$, then $a\ne\, 1,2,3$,
 since $8+1=9,8+2=10, 8+3=11$ are 
  not in $E.$  Similarly,
   $a\ne 5,6,7,9,10,11$
 since none of those elements is in $E$. Thus, $a\in \{0,4,8\}$ and the first set is seen not be a sum.
 
 Of course, $\{0,4,8\}+\{0,1,2\}=\{0,1,2,4,5,6,8,9,10\}$ and both final conclusions follow.

(1). For the extremality of the first set, let $\mu =
 \delta_{0}  +   e^{ 23\pi i  / 12}\delta_{1}  +   e^{ 3\pi i  / 2}\delta_{2}  +   e^{ 17\pi i  / 12}\delta_{3}  +  \delta_{4}  +   e^{\pi i / 4}\delta_{5}  +   e^{ 7\pi i  / 6}\delta_{6}  +   e^{ 5\pi i  / 12}\delta_{7}  +   e^{ 4\pi i  / 3}\delta_{8}
.$
  Then
\begin{align*}\mu*\tilde\mu &= 9\delta_0+ \big( e^{ 23 \pi i / 12}+ e^{ 19 \pi i / 12}+ e^{ 23 \pi i / 12}+ e^{ 7 \pi i / 12}+ e^{\pi i / 4}
\\&\quad\qquad+ e^{ 11 \pi i / 12}+ e^{ 5 \pi i / 4}+ e^{ 11 \pi i / 12}\big)\delta_{1} 
\\&\qquad  +  \big( -i - i+  i + e^{ 5 \pi i / 6}+ e^{ 7 \pi i / 6}+ e^{\pi i / 6}+ e^{\pi i / 6}\big)\delta_{2} 
\\&\qquad  +  \big( e^{ 17 \pi i / 12}+ e^{\pi i / 12}+ e^{ 3 \pi i / 4}+ e^{ 7 \pi i / 4}+ e^{ 5 \pi i / 12}+ e^{ 13 \pi i / 12}\big)\delta_{3}
 \\&\qquad  +  \big( e^{ 2 \pi i / 3}+ 1+ e^{\pi i / 3}+ e^{ 5 \pi i / 3}+ e^{ 1 \pi i}+ e^{ 4 \pi i / 3}\big)\delta_{4} 
\\&\qquad  +  \big( e^{ 19 \pi i / 12}+ e^{ 7 \pi i / 12}+ e^{\pi i / 4}+ e^{ 5 \pi i / 4}+ e^{ 11 \pi i / 12}+ e^{ 23 \pi i / 12}\big)\delta_{5} 
\\&\qquad  +  \big( e^{ 5 \pi i / 6}- i+ e^{\pi i / 6}+ e^{ 7 \pi i / 6}+  i + e^{ 11 \pi i / 6}\big)\delta_{6} 
\\&\qquad  +  \big( e^{ 7 \pi i / 4}+ e^{ 3 \pi i / 4}+ e^{ 13 \pi i / 12}+ e^{\pi i / 12}+ e^{ 5 \pi i / 12}+ e^{ 17 \pi i / 12}\big)\delta_{7} 
\\&\qquad  +  \big( 1+ e^{ 5 \pi i / 3}+ e^{\pi i / 3}+ e^{ 1 \pi i}+ e^{ 2 \pi i / 3}+ e^{ 4 \pi i / 3}\big)\delta_{8} 
\\&\qquad  +  \big( e^{ 7 \pi i / 12}+ e^{ 23 \pi i / 12}+ e^{ 5 \pi i / 4}+ e^{\pi i / 4}+ e^{ 19 \pi i / 12}+ e^{ 11 \pi i / 12}\big)\delta_{9} 
\\&\qquad  +  \big(  i +  i - i+ e^{ 7 \pi i / 6}+ e^{ 5 \pi i / 6}+ e^{ 11 \pi i / 6}+ e^{ 11 \pi i / 6}\big)\delta_{10} 
\\&\qquad  +  \big( e^{\pi i / 12}+ e^{ 5 \pi i / 12}+ e^{\pi i / 12}+ e^{ 17 \pi i / 12}+ e^{ 7 \pi i / 4}
\\&\quad\qquad+ e^{ 13 \pi i / 12}+ e^{ 3 \pi i / 4}+ e^{ 13 \pi i}\big)\delta_{11} 
\\&=9\delta_0.
\end{align*}

(2).  Let $  
\mu = \delta_{0}  +   e^{ 10\pi i  / 6}\delta_{1}  +   e^{ 10\pi i  / 6}\delta_{2}  +   e^{ 8\pi i  / 6}\delta_{4}  +   e^{ 102\pi i  / 6}\delta_{5}  +   e^{ 2\pi i  / 6}\delta_{6}  +   e^{ 4\pi i  / 6}\delta_{8}  +   e^{ 10\pi i  / 6}\delta_{9}  +   e^{ 6\pi i  / 6}\delta_{10}
.$
  Then
\begin{align*}\mu * \tilde\mu&=    9\delta_{0}   +  \big( e^{ 5 \pi i / 3}+ 1+ e^{ 5 \pi i / 3}+ e^{ 4 \pi i / 3}+ e^{\pi  i}+ e^{ 4 \pi i / 3}\big)\delta_{1} \\&\qquad  +  \big( e^{\pi  i}+ e^{ 5 \pi i / 3}+ e^{ 5 \pi i / 3}+ e^{\pi  i}+ e^{\pi i / 3}+ e^{\pi i / 3}\big)\delta_{2} 
\\&\qquad  +  \big( e^{\pi i / 3}+ e^{ 2 \pi i / 3}+ e^{ 5 \pi i / 3}+ e^{ 4 \pi i / 3}+ e^{ 5 \pi i / 3}+ e^{ 4 \pi i / 3}\big)\delta_{3} \\&\qquad  +  \big( e^{ 4 \pi i / 3}+ 1+ e^{ 2 \pi i / 3}+ e^{ 4 \pi i / 3}+ e^{ 4 \pi i / 3}+ e^{ 2 \pi i / 3}+ e^{ 4 \pi i / 3}
\\& \quad\qquad + e^{ 2 \pi i / 3}+ e^{ 2 \pi i / 3}\big)\delta_{4} 
\\&\qquad  +  \big( e^{\pi  i}+ 1+ e^{\pi  i}+ e^{ 2 \pi i / 3}+ e^{\pi i / 3}+ 1\big)\delta_{5} \\&\qquad  +  \big( e^{ 5 \pi i / 3}+ e^{\pi  i}+ e^{\pi i / 3}+ e^{\pi i / 3}+ e^{\pi  i}+ e^{ 5 \pi i / 3}\big)\delta_{6} 
\\&\qquad  +  \big( e^{\pi  i}+ e^{ 4 \pi i / 3}+ e^{ 5 \pi i / 3}+ 1+ e^{\pi  i}+ 1\big)\delta_{7} \\&\qquad  +  \big( e^{ 2 \pi i / 3}+ e^{ 2 \pi i / 3}+ e^{ 4 \pi i / 3}+ e^{ 2 \pi i / 3}+ e^{ 4 \pi i / 3}\\& \quad\qquad + e^{ 4 \pi i / 3}+ e^{ 2 \pi i / 3}+ 1+ e^{ 4 \pi i / 3}\big)\delta_{8} 
\end{align*}
\begin{align*}
\phantom{\mu * \tilde\mu}&\qquad  +  \big( e^{\pi i / 3}+ e^{ 2 \pi i / 3}+ e^{\pi i / 3}+ e^{ 2 \pi i / 3}+ e^{ 5 \pi i / 3}+ e^{ 4 \pi i / 3}\big)\delta_{9} \\&\qquad  +  \big( e^{\pi i / 3}+ e^{\pi i / 3}+ e^{\pi  i}+ e^{ 5 \pi i / 3}+ e^{ 5 \pi i / 3}+ e^{\pi  i}\big)\delta_{10} \\&\qquad  +  \big( e^{\pi i / 3}+ 1+ e^{\pi i / 3}+ e^{ 2 \pi i / 3}+ e^{\pi  i}+ e^{ 2 \pi i}\big)\delta_{11} 
   \\&=9\delta_0. 
\end{align*}
\end{proof}


\begin{prop}\label{prop9in13} The 9-element set
$\{0, 1, 2, 3, 4, 5, 7, 9, 10\}$ is extreme in \Zp{13}{}.
\end{prop}

\begin{proof}
Let $\mu = \delta_{0} +\delta_{1} - \delta_{2} - \delta_{3} +\delta_{4} - \delta_{5} - \delta_{7} - \delta_{9} - \delta_{10} 
$.  Then 
\begin{align*} 
\mu * \tilde \mu & = 9\delta_{0}   +  \big(1-1+1-1- 1 + 1\big)\delta_{1} 
%\\&\qquad 
 +  \big(  -1 - 1 - 1 + 1+ 1+ 1\big)\delta_{2} \\&\qquad  +  \big(  -1 - 1 + 1+1-1+ 1\big)\delta_{3} 
 %\\&\qquad 
  +  \big(  -1 - 1 +1-1+ 1+ 1\big)\delta_{4} \\&\qquad  +  \big(  -1 +1-1+1-1+ 1\big)\delta_{5} 
 % \\&\qquad 
   +  \big(  -1 + 1+1-1+1-1\big)\delta_{6} \\&\qquad  +  \big(  -1 +1-1- 1 + 1+ 1\big)\delta_{7} 
%   \\&\qquad 
    +  \big(  -1 +1-1+1-1+ 1\big)\delta_{8} \\&\qquad  +  \big(1-1+ 1+1-1- 1 \big)\delta_{9} 
 %  \\&\qquad
     +  \big(  -1 + 1+1-1+1-1\big)\delta_{10} \\&\qquad  +  \big(  -1 - 1 - 1 + 1+ 1+ 1\big)\delta_{11} \
   %\&\qquad 
    +  \big(1-1+1-1- 1 + 1\big)\delta_{12} 
\\&=9\delta_0.\qedhere
\end{align*}
\end{proof}
\begin{comment}
\begin{prop}
$\{0,1,3,4,6,7,9,10\}$
is extreme in \Zp{12}{}.
\end{prop}
\begin{proof}
Let $\mu =
 \delta_{0}  +   e^{ 7\pi i  / 6}\delta_{1}  +  \delta_{3}  +   e^{ 5\pi i  / 3}\delta_{4}  +   e^{\pi i }\delta_{6}  +   e^{ 7\pi i  / 6}\delta_{7}  +   e^{\pi i }\delta_{9}  +   e^{ 5\pi i  / 3}\delta_{10}.  
$
Then
\begin{align*}   \mu*\tilde\mu =8\delta_{0}&  +  \big( e^{ 7 \pi i / 6}+ e^{ 5 \pi i / 3}+ e^{\pi i / 6}+ e^{ 2 \pi i / 3}\big)\delta_{1} 
  +  \big( e^{\pi i / 3}+ e^{ 5 \pi i / 6}+ e^{ 4 \pi i / 3}+ e^{ 11 \pi i / 6}\big)\delta_{2} 
  \\&  +  \big( e^{\pi i }- i+ 1+  i + e^{\pi i }- i+ 1+  i \big)\delta_{3}
   \\&  +  \big( e^{\pi i / 6}+ e^{ 5 \pi i / 3}+ e^{ 7 \pi i / 6}+ e^{ 2 \pi i / 3}\big)\delta_{4} 
    +  \big( e^{ 5 \pi i / 6}+ e^{\pi i / 3}+ e^{ 11 \pi i / 6}+ e^{ 4 \pi i / 3}\big)\delta_{5} 
   \\&  +  \big( e^{\pi i }+ 1+ e^{\pi i }+ 1+ e^{\pi i }+ 1+ e^{\pi i }+ 1\big)\delta_{6} 
   \\&  +  \big( e^{\pi i / 6}+ e^{ 2 \pi i / 3}+ e^{ 7 \pi i / 6}+ e^{ 5 \pi i / 3}\big)\delta_{7} 
    +  \big( e^{\pi i / 3}+ e^{ 5 \pi i / 6}+ e^{ 4 \pi i / 3}+ e^{ 11 \pi i / 6}\big)\delta_{8} 
   \\&  +  \big( 1- i+ e^{\pi i }+  i + 1- i+ e^{\pi i }+  i \big)\delta_{9} 
    +  \big( e^{ 7 \pi i / 6}+ e^{ 2 \pi i / 3}+ e^{\pi i / 6}+ e^{ 5 \pi i / 3}\big)\delta_{10} 
   \\&  +  \big( e^{ 5 \pi i / 6}+ e^{\pi i / 3}+ e^{ 11 \pi i / 6}+ e^{ 4 \pi i }\big)\delta_{11} = 8\delta_0.
   \qedhere
\end{align*}
\end{proof}
\end{comment}


\begin{comment}% This set is the sum of a subgroup and an extreme set, {0,2,4,8}.
\begin{prop}
$\{0,1,2,4,7,8,9,11\}$ is extreme in \Zp{14}{}.
\end{prop}

\begin{proof}
Let $\mu =
  \delta_{0}  +    i \delta_{1}  +    -\delta_{2}  +    -\delta_{4}  +   e^{ 3\pi i  / 2}\delta_{7}  +    -\delta_{8}  +    i \delta_{9}  +    i \delta_{11}$. Then
\begin{align*} \mu * \tilde\mu=
 8\delta_{0} &  +  \big(  i +  i - i- i\big)\delta_{1} 
  +  \big(  -1 +1-1+ 1\big)\delta_{2} 
 \\&  +  \big( -i +  i +  i - i\big)\delta_{3} 
  +  \big(1-1+1-1\big)\delta_{4} \\&  +  \big( -i +  i +  i - i\big)\delta_{5} 
    +  \big(  -1 +1-1+ 1\big)\delta_{6} \\&  +  \big(  i - i+  i +  i - i+  i - i- i\big)\delta_{7} \\&  +  \big(  -1 +1-1+ 1\big)\delta_{8} 
      +  \big( -i +  i +  i - i\big)\delta_{9} \\&  +  \big(  -1 +1-1+ 1\big)\delta_{10} 
       +  \big( -i - i+  i +  i \big)\delta_{11} \\&  +  \big(  -1 +1-1+ 1\big)\delta_{12} 
        +  \big( -i - i+  i - 1 \big)\delta_{13} \\& = 8\delta_0.
        \qedhere
\end{align*}
\end{proof}
\end{comment}

\subsection{Sets with 10 elements}
The only 10-element extreme sets we have found are sums of 5-element cosets with 2-element cosets and 5-element seubsets of
6-element cosets with 2-element cosets.


 \subsection{Sets with 11 elements}
 
\begin{prop}\label{prop11inZ12}
The 11-element set $\{0,  1,  2,  3,  4,  5,  6,  7,  8,  9,  10 \}$ is extreme in $\Zp{12}{}$. 
\newline \end{prop} 


 
\begin{proof}
Let  $\mu =  \delta_{0}+e^{7\pii/4}\delta_{1}-i\delta_{2}+e^{5\pii/4}\delta_{3}+\delta_{4}+e^{7\pii/4}\delta_{5}+i\delta_{6}+e^{\pii/4}\delta_{7}-\delta_{8}+e^{7\pii/4}\delta_{9}+i\delta_{10}
$.  Then 
\begin{align*} 
\mu * \tilde \mu & = \big(1+1+1+1+1+1+1+1+1+1+1\big)\delta_{0}
\\&\qquad
+\big(e^{7\pii/4}+e^{7\pii/4}+e^{7\pii/4}+e^{3\pii/4}+e^{7\pii/4}+e^{3\pii/4}
\\&\qquad\qquad +e^{7\pii/4}+e^{3\pii/4}+e^{3\pii/4}+e^{3\pii/4}-i-i-i+i+
\\&\qquad\qquad
i+i+i+i-i-i+e^{\pii/4}
+e^{5\pii/4}+e^{5\pii/4}+e^{\pii/4}+e^{\pii/4}+e^{5\pii/4}
\\&\qquad\qquad
+e^{\pii/4}+e^{5\pii/4}+e^{5\pii/4}
\\&\qquad\qquad
+e^{\pii/4}-1+1-1+1+1-1-1-1+1+1
\\&\qquad\qquad
+e^{7\pii/4}+e^{3\pii/4}+e^{7\pii/4}+e^{3\pii/4}+e^{7\pii/4}+e^{3\pii/4}+e^{3\pii/4}
\\&\qquad\qquad+e^{7\pii/4}+e^{7\pii/4}+e^{3\pii/4}-i-i+i-i-i+i+i-i+i+i\big)\delta_{6}
\\&\qquad
+\big(e^{\pii/4}+e^{5\pii/4}+e^{5\pii/4}+e^{\pii/4}+e^{\pii/4}+e^{5\pii/4}
\\&\qquad\qquad
+e^{\pii/4}+e^{5\pii/4}+e^{\pii/4}+e^{5\pii/4}\big)\delta_{7}
\\&\qquad
+\big(1+1-1-1-1+1+1-1+1-1\big)\delta_{8}
\\&\qquad
+\big(e^{3\pii/4}+e^{7\pii/4}+e^{7\pii/4}+e^{3\pii/4}+e^{7\pii/4}+e^{3\pii/4}+e^{3\pii/4}
\\&\qquad\qquad+e^{7\pii/4}+e^{7\pii/4}+e^{3\pii/4}\big)\delta_{9}
\\&\qquad
+\big(i+i-i-i-i-i-i+i+i+i\big)\delta_{10}
\\&\qquad
+\big(e^{\pii/4}+e^{\pii/4}+e^{\pii/4}+e^{5\pii/4}+e^{\pii/4}+e^{5\pii/4}+e^{\pii/4}
\\&\qquad\qquad+e^{5\pii/4}+e^{5\pii/4}+e^{5\pii/4}\big)\delta_{11}\\&=11\delta_{0 }\qedhere
\end{align*}
\end{proof}



\subsection{Sets with 12 elements}
 
The only 12-element extreme sets we have found are sums of 6-element cosets with 2-element cosets.

 \subsection{Sets with 13 elements}
 
 

%13inZ14
$ \delta_{0}+\delta_{12}+e^{4\pii/3}(\delta_{1}+ \delta_{3} + \delta_{4} +\delta_{8}+ \delta_{9}+\delta_{11})
+e^{2\pii/3}(\delta_{2}+ \delta_{5}+ \delta_{6}+ \delta_{7}+ \delta_{10})$

Let  $\mu =  \delta_{(0,0)}+e^{4\pii/3}\delta_{(0,1)}+e^{2\pii/3}\delta_{(0,2)}+e^{4\pii/3}\delta_{(0,3)}+e^{4\pii/3}\delta_{(0,4)}+e^{2\pii/3}\delta_{(0,5)}+e^{2\pii/3}\delta_{(0,6)}+e^{2\pii/3}\delta_{(0,7)}+e^{4\pii/3}\delta_{(0,8)}+e^{4\pii/3}\delta_{(0,9)}+e^{2\pii/3}\delta_{(0,10)}+e^{4\pii/3}\delta_{(0,11)}+\delta_{(0,12)}
$.  Then \begin{align*} 
\mu * \tilde \mu & = \big(1+1+1+1+1+1+1+1+1+1+1+1+1\big)\delta_{(0,0)}
\\&\qquad
+\big(e^{4\pii/3}+e^{4\pii/3}+e^{2\pii/3}+1+e^{4\pii/3}+1+1+e^{2\pii/3}+1+e^{4\pii/3}+e^{2\pii/3}+e^{2\pii/3}\big)\delta_{(0,1)}
\\&\qquad
+\big(1+e^{2\pii/3}+1+e^{2\pii/3}+e^{4\pii/3}+e^{4\pii/3}+1+e^{2\pii/3}+e^{2\pii/3}+e^{4\pii/3}+1+e^{4\pii/3}\big)\delta_{(0,2)}
\\&\qquad
+\big(e^{2\pii/3}+e^{4\pii/3}+e^{4\pii/3}+1+1+e^{4\pii/3}+e^{4\pii/3}+e^{2\pii/3}+e^{2\pii/3}+1+1+e^{2\pii/3}\big)\delta_{(0,3)}
\\&\qquad
+\big(e^{4\pii/3}+1+e^{2\pii/3}+e^{4\pii/3}+e^{4\pii/3}+1+e^{4\pii/3}+1+e^{2\pii/3}+1+e^{2\pii/3}+e^{2\pii/3}\big)\delta_{(0,4)}
\\&\qquad
+\big(e^{2\pii/3}+e^{2\pii/3}+e^{4\pii/3}+e^{4\pii/3}+e^{2\pii/3}+e^{4\pii/3}+1+1+1+1+e^{2\pii/3}+e^{4\pii/3}\big)\delta_{(0,5)}
\\&\qquad
+\big(e^{2\pii/3}+1+1+1+e^{4\pii/3}+e^{2\pii/3}+e^{4\pii/3}+e^{2\pii/3}+1+e^{4\pii/3}+e^{2\pii/3}+e^{4\pii/3}\big)\delta_{(0,6)}
\\&\qquad
+\big(e^{4\pii/3}+1+e^{4\pii/3}+e^{2\pii/3}+1+e^{2\pii/3}+e^{2\pii/3}+1+e^{2\pii/3}+e^{4\pii/3}+1+e^{4\pii/3}\big)\delta_{(0,7)}
\\&\qquad
+\big(e^{4\pii/3}+e^{2\pii/3}+e^{4\pii/3}+1+e^{2\pii/3}+e^{4\pii/3}+e^{2\pii/3}+e^{4\pii/3}+1+1+1+e^{2\pii/3}\big)\delta_{(0,8)}
\\&\qquad
+\big(e^{4\pii/3}+e^{2\pii/3}+1+1+1+1+e^{4\pii/3}+e^{2\pii/3}+e^{4\pii/3}+e^{4\pii/3}+e^{2\pii/3}+e^{2\pii/3}\big)\delta_{(0,9)}
\\&\qquad
+\big(e^{2\pii/3}+e^{2\pii/3}+1+e^{2\pii/3}+1+e^{4\pii/3}+1+e^{4\pii/3}+e^{4\pii/3}+e^{2\pii/3}+1+e^{4\pii/3}\big)\delta_{(0,10)}
\\&\qquad
+\big(e^{2\pii/3}+1+1+e^{2\pii/3}+e^{2\pii/3}+e^{4\pii/3}+e^{4\pii/3}+1+1+e^{4\pii/3}+e^{4\pii/3}+e^{2\pii/3}\big)\delta_{(0,11)}
\\&\qquad
+\big(e^{4\pii/3}+1+e^{4\pii/3}+e^{2\pii/3}+e^{2\pii/3}+1+e^{4\pii/3}+e^{4\pii/3}+e^{2\pii/3}+1+e^{2\pii/3}+1\big)\delta_{(0,12)}
\\&\qquad
+\big(e^{2\pii/3}+e^{2\pii/3}+e^{4\pii/3}+1+e^{2\pii/3}+1+1+e^{4\pii/3}+1+e^{2\pii/3}+e^{4\pii/3}+e^{4\pii/3}\big)\delta_{(0,13)}\\&=13\delta_{(0 , 0)}\qedhere
\end{align*}


 \begin{prop}\label{prop13in14}
  $\{0-12\}$ is extreme in \Zp{14}{}.
 \end{prop}
 xxx
 \begin{proof} Let %\begin{align*} 
$\mu =
%&
  \delta_{0}  +   e^{ 7\pi i  / 4}\delta_{1}  +   \delta_{2}  +   e^{ 3\pi i  / 4}\delta_{3}  +   e^{ 3\pi i  / 2}\delta_{4}  +   e^{ 5\pi i  / 4}\delta_{5}  %\\&
+   i\delta_{6}  +   e^{ \pi i  / 4}\delta_{7}  +   e^{ 3\pi i  / 4}\delta_{8}  +   e^{ 7\pi i  / 4}\delta_{9}  +  \delta_{10}  +   e^{ 3\pi i  / 4}\delta_{11}  +   \delta_{12}.
$%\end{align*}

  Then
\begin{align*}\mu * \tilde\mu&=   13\delta_{0}   +  \big( e^{ 7 \pi i / 4}+ e^{ 9 \pi i / 4}+ e^{ 3 \pi i / 4}+ e^{ 11 \pi i / 4}+ e^{ 7 \pi i / 4}+ e^{ 7 \pi i / 4}
\\& \quad\qquad + e^{ 5 \pi i / 4}+ e^{ 13 \pi i / 4}+ e^{ 9 \pi i / 4}+ e^{ 9 \pi i / 4}+ e^{ 3 \pi i / 4}+ e^{ 13 \pi i / 4}\big)\delta_{1} 
\\& \qquad  +  \big( 1+ 1+ e^{ \pi  i}- i+ e^{ 5 \pi i / 2}- i+ e^{ \pi  i}
\\& \quad\qquad + e^{ 5 \pi i / 2}+ e^{ 7 \pi i / 2}+ e^{ 5 \pi i / 2}+ e^{ \pi  i}+ 1\big)\delta_{2} 
\\& \qquad  +  \big( e^{ 5 \pi i / 4}+ e^{ 7 \pi i / 4}+ e^{ 3 \pi i / 4}+ e^{ 7 \pi i / 4}+ e^{ 5 \pi i / 4}+ e^{ 9 \pi i / 4}+ e^{ 3 \pi i / 4}
\\& \quad\qquad+ e^{ 9 \pi i / 4}+ e^{ 11 \pi i / 4}+ e^{ 15 \pi i / 4}+ e^{ 5 \pi i / 4}+ e^{ 9 \pi i / 4}\big)\delta_{3} 
\\& \qquad  +  \big( 1+ e^{ 3 \pi i}+ 1- i- i+ e^{ \pi  i}- i+ 1
\\& \quad\qquad
+ e^{ 5 \pi i / 2}+ e^{ 3 \pi i}+ e^{ 5 \pi i / 2}+ e^{ 5 \pi i / 2}\big)\delta_{4} 
\\& \qquad  +  \big( e^{\pi i / 4}+ e^{ 7 \pi i / 4}+ e^{ 13 \pi i / 4}+ e^{ 3 \pi i / 4}+ e^{ 5 \pi i / 4}+ e^{ 5 \pi i / 4}+ e^{\pi i / 4}
\\& \quad\qquad
+ e^{ 11 \pi i / 4}+ e^{ 7 \pi i / 4}+ e^{ 15 \pi i / 4}\big)\delta_{5} 
\\& \qquad  +  \big(  i + 1+ 1+ 1- i+ e^{ \pi  i}+  i 
\\& \quad\qquad
- i+ e^{ 3 \pi i}+ e^{ 5 \pi i / 2}- i+ e^{ 3 \pi i}\big)\delta_{6}
 \\& \qquad  +  \big( e^{ 7 \pi i / 4}+ e^{ 9 \pi i / 4}+ e^{ 9 \pi i / 4}+ e^{ 3 \pi i / 4}+ e^{ 11 \pi i / 4}+ e^{ 5 \pi i / 4}
\\& \quad\qquad
+ e^{\pi i / 4}+ e^{ 7 \pi i / 4}+ e^{ 7 \pi i / 4}+ e^{ 13 \pi i / 4}+ e^{ 5 \pi i / 4}+ e^{ 11 \pi i / 4}\big)\delta_{7} 
 \\& \qquad  +  \big( e^{ \pi  i}+ e^{ 7 \pi i / 2}+ e^{ 5 \pi i / 2}+ e^{ \pi  i}- i
\\& \quad\qquad
+ e^{ 5 \pi i / 2}+ e^{ \pi  i}- i+ 1+ 1+ 1+ e^{ 5 \pi i / 2}\big)\delta_{8} 
 \\& \qquad  +  \big( e^{ 3 \pi i / 4}+ e^{ 11 \pi i / 4}+ e^{ 15 \pi i / 4}+ e^{ 5 \pi i / 4}+ e^{ 7 \pi i / 4}+ e^{ 5 \pi i / 4}
\\& \quad\qquad
+ e^{ 9 \pi i / 4}+ e^{\pi i / 4}+ e^{ 7 \pi i / 4}+ e^{ 9 \pi i / 4}+ e^{ 3 \pi i / 4}+ e^{ 13 \pi i / 4}\big)\delta_{9} 
 \\& \qquad  +  \big(  i + e^{ 5 \pi i / 2}+ e^{ 3 \pi i}+ e^{ 5 \pi i / 2}+ 1- i+ e^{ \pi  i}
\\& \quad\qquad
- i- i+ 1+ e^{ \pi  i}+ 1\big)\delta_{10} 
 \\& \qquad  +  \big( e^{ 5 \pi i / 4}+ e^{ 9 \pi i / 4}+ e^{ 11 \pi i / 4}+ e^{ 7 \pi i / 4}+ e^{ 13 \pi i / 4}+ e^{ 7 \pi i / 4}
\\& \quad\qquad
+ e^{ 5 \pi i / 4}+ e^{\pi i / 4}+ e^{ 11 \pi i / 4}+ e^{ 7 \pi i / 4}+ e^{ 3 \pi i / 4}+ e^{ 9 \pi i / 4}\big)\delta_{11} 
 \\& \qquad  +  \big( 1+ e^{ 3 \pi i}+ e^{ 5 \pi i / 2}- i+ e^{ 5 \pi i / 2}
\\& \quad\qquad
+ e^{ 3 \pi i}- i+  i - i+ e^{ 3 \pi i}+ 1+ 1\big)\delta_{12} 
 \\& \qquad  +  \big( e^{\pi i / 4}+ e^{ 7 \pi i / 4}+ e^{ 13 \pi i / 4}+ e^{ 5 \pi i / 4}+ e^{ 9 \pi i / 4}
+ e^{ 9 \pi i / 4}+ e^{ 11 \pi i / 4}\\& \quad\qquad+ e^{ 3 \pi i / 4}+ e^{ 7 \pi i / 4}+ e^{ 7 \pi i / 4}+ e^{ 13 \pi i / 4}+ e^{ 3 \pi i}\big)\delta_{13} \
\\&= 13\delta_0.
\qedhere
\end{align*}
\end{proof}

\subsection{Sets with 14 to 16 elements}
Except for the special case in the next subsection, the computatational time needed to search for extreme
sets with more than 13 elements is prohibitive at this writing.

\begin{prop}\label{proppp16inZ16}
The 16 element set $\{0,  1,  2,  3,  4,  5,  6,  7,  8,  9,  10,  11,  12,  13,  14,  15 \}$ is extreme in $\Zp{16}{}$. \newline \end{prop} 


 
\begin{proof}
Let  $\mu =  \delta_{0}+\delta_{1}+i\delta_{2}+\delta_{3}+\delta_{4}+i\delta_{5}+\delta_{6}-\delta_{7}+\delta_{8}-\delta_{9}-i\delta_{10}+\delta_{11}+\delta_{12}-i\delta_{13}-\delta_{14}+\delta_{15}
$.  

Then 
\begin{align*} 
\mu * \tilde \mu & = \big(1+1+1+1+1+1+1+1+1+1+1+1+1+1+1+1\big)\delta_{0}
\\&\qquad
+\big(1+1+i-i+1+i-i-1-1-1+i+i+1-i-i-1\big)\delta_{1}
\\&\qquad
+\big(-1+1+i+1-i+i+1+i+1+1-i-1+i-i-1+i\big)\delta_{2}
\\&\qquad
+\big(i-1+i+1+1+1+1-1-i-1+i+1-1+1-1+1\big)\delta_{3}
\\&\qquad
+\big(1+i-i+1+1+i-i-1+1+i-i-1+1+i-i+1\big)\delta_{4}
\\&\qquad
+\big(1+1-1-1+1+i+1+i+1-1-1+1-1-i+1+i\big)\delta_{5}
\\&\qquad
+\big(i+1+i+i-1+i+1-1-i-1-i-i+1+i-1-1\big)\delta_{6}
\\&\qquad
+\big(-1+i+i+1+i-i+1-1+1+i-i+1-i-i+1+1\big)\delta_{7}
\\&\qquad
+\big(1-1-1+1+1-1-1-1+1-1-1+1+1-1-1-1\big)\delta_{8}
\\&\qquad
+\big(-1+1-i+i+1+i+i+1+1-1-i-i+1-i+i+1\big)\delta_{9}
\\&\qquad
+\big(1-1+i-1+i+i+1-i-1-1-i+1-i-i-1-i\big)\delta_{10}
\\&\qquad
+\big(-i+1-i+1-1-1+1-1+i+1-i+1+1-1-1+1\big)\delta_{11}
\\&\qquad
+\big(1-i+i-1+1-i+i-1+1-i+i+1+1-i+i+1\big)\delta_{12}
\\&\qquad
+\big(1+1+1+1-1+i-1-i+1-1+1-1+1-i-1-i\big)\delta_{13}
\\&\qquad
+\big(-i+1+i-i+1-i+1+1+i-1-i+i-1-i-1+1\big)\delta_{14}
\\&\qquad
+\big(1-i+i+1-i+i-1-1-1-i-i+1+i+i-1+1\big)\delta_{15}\\&=16\delta_{0 }\qedhere
\end{align*}
\end{proof}



\begin{prop}\label{prop16inZ17}
The 16 element set $\{0,  1,  2,  3,  4,  5,  6,  7,  8,  9,  10,  11,  12,  13,  14,  15 \}$ is extreme in $ \Zp{17}{}$. \newline \end{prop} 


 
\begin{proof}
Let  $\mu =  \delta_{0}+e^{6\pii/5}\delta_{1}+e^{6\pii/5}\delta_{2}+e^{8\pii/5}\delta_{3}+e^{2\pii/5}\delta_{4}+e^{8\pii/5}\delta_{5}+\delta_{6}+e^{2\pii/5}\delta_{7}+e^{2\pii/5}\delta_{8}+\delta_{9}+e^{8\pii/5}\delta_{10}+e^{2\pii/5}\delta_{11}+e^{8\pii/5}\delta_{12}+e^{6\pii/5}\delta_{13}+e^{6\pii/5}\delta_{14}+\delta_{15}
$.  Then 
\begin{align*} 
\mu * \tilde \mu & = 16\delta_{0}
%\\&\quad
+\big(3e^{6\pii/5}+3+3e^{2\pii/5}+3e^{4\pii/5}+3e^{8\pii/5}\big)
%\\ &\qquad\quad\times
\big(\delta_{1}
+\delta_{2}+\cdots+\delta_{16}\big)\\&=16\delta_{0 }\qedhere
\end{align*}
\end{proof}

\begin{prop}\label{prop16inZ16}
The 16 element set $\{0,  1,  2,  3,  4,  5,  6,  7,  8,  9,  10,  11,  12,  13,  14,  15 \}$ is extreme in $ \times  \times  \times  \times  \times  \times  \times  \times  \times  \times  \times  \times  \times \Zp{16}{}$. \newline \end{prop} 


 
 
\begin{proof}
Let  $\mu =  \delta_{0}+\delta_{1}+\delta_{2}-i\delta_{3}-\delta_{4}+i\delta_{5}+\delta_{6}-\delta_{7}+\delta_{8}-\delta_{9}+\delta_{10}-i\delta_{11}-\delta_{12}+i\delta_{13}+\delta_{14}+\delta_{15}
$.  Then 
\begin{align*} 
\mu * \tilde \mu & = \big(1+1+1+1+1+1+1+1+1+1+1+1+1+1+1+1\big)\delta_{0}
\\&\qquad
+\big(1+1+1-i-i-i-i-1-1-1-1-i-i-i-i+1\big)\delta_{1}
\\&\qquad
+\big(1+1+1-i-1-1-1+i+1+1+1+i-1-1-1-i\big)\delta_{2}
\\&\qquad
+\big(-i+1+1-i-1+i+i+1-i-1-1-i+1+i+i-1\big)\delta_{3}
\\&\qquad
+\big(-1-i+1-i-1+i+1-i-1+i+1+i-1-i+1+i\big)\delta_{4}
\\&\qquad
+\big(i-1-i-i-1+i+1-1+i+1-i-i+1+i-1+1\big)\delta_{5}
\\&\qquad
+\big(1+i-1-1-1+i+1-1+1-i-1-1-1-i+1-1\big)\delta_{6}
\\&\qquad
+\big(-1+1+i+i+i+i+1-1+1-1+i+i+i+i-1+1\big)\delta_{7}
\\&\qquad
+\big(1-1+1+1+1+1+1-1+1-1+1+1+1+1+1-1\big)\delta_{8}
\\&\qquad
+\big(-1+1-1-i-i-i-i-1+1-1+1-i-i-i-i+1\big)\delta_{9}
\\&\qquad
+\big(1-1+1+i-1-1-1+i+1-1+1-i-1-1-1-i\big)\delta_{10}
\\&\qquad
+\big(-i+1-1-i+1+i+i+1-i-1+1-i-1+i+i-1\big)\delta_{11}
\\&\qquad
+\big(-1-i+1+i-1-i+1-i-1+i+1-i-1+i+1+i\big)\delta_{12}
\\&\qquad
+\big(i-1-i-i+1+i-1-1+i+1-i-i-1+i+1+1\big)\delta_{13}
\\&\qquad
+\big(1+i-1-1-1-i+1+1+1-i-1-1-1+i+1+1\big)\delta_{14}
\\&\qquad
+\big(1+1+i+i+i+i-1-1-1-1+i+i+i+i+1+1\big)\delta_{15}\\&=16\delta_{0 }\qedhere
\end{align*}
\end{proof}





 
\subsection{Sets with 17 elements}

\begin{prop}\label{prop17in18}
$\{0, 1, 2, 3, 4, 5, 6, 7, 8, 9, 10, 11, 12, 13, 14, 15, 16\}$ is extreme in \Zp{18}{}.
\end{prop}
\begin{proof}
Let  $\mu =  \delta_{(0,0)}-i\delta_{(1,0)}-\delta_{(2,0)}-i\delta_{(3,0)}+\delta_{(4,0)}-i\delta_{(5,0)}-\delta_{(6,0)}+i\delta_{(7,0)}-\delta_{(8,0)}+i\delta_{(9,0)}-\delta_{(10,0)}-i\delta_{(11,0)}+\delta_{(12,0)}-i\delta_{(13,0)}-\delta_{(14,0)}-i\delta_{(15,0)}+\delta_{(16,0)}
$.  Then 
\begin{align*} 
\mu * \tilde \mu & = \big(1+1+1+1+1+1+1+1+1+1+1+1+1+1+1+1+1\big)\delta_{(0,0)}
\\&\qquad
+\big(-i-i+i+i-i-i-i+i-i+i+i+i-i-i+i+i
\\&\qquad\qquad+1-1+1-1+1-1-1+1+1+1-1-1+1-1+1-1
\\&\qquad\qquad+i-i-i+i+i-i+i-i-i+i+i-i+i-i-i+i
\\&\qquad\qquad-1+1-1+1+1+1-1-1-1+1-1-1-1+1+1+1
\\&\qquad\qquad+i+i-i-i-i-i-i-i+i-i+i-i+i+i+i+i
\\&\qquad\qquad+1+1+1+1+1-1-1+1-1-1+1-1-1+1-1-1
\\&\qquad\qquad+i-i-i+i+i-i+i-i-i-i-i+i+i+i+i-i
\\&\qquad\qquad-1+1-1+1-1+1-1-1-1+1+1+1+1+1-1-1
\\&\qquad\qquad-i+i-i-i+i+i-i+i+i-i+i+i-i-i+i-i
\\&\qquad\qquad-1-1+1+1+1+1+1-1-1-1+1-1+1-1+1-1
\\&\qquad\qquad-i+i+i+i+i-i-i-i-i+i-i+i+i-i-i+i
\\&\qquad\qquad-1-1+1-1-1+1-1-1+1-1-1+1+1+1+1+1\big)\delta_{(12,0)}
\\&\qquad
+\big(i+i+i+i-i+i-i+i-i-i-i-i-i-i+i+i\big)\delta_{(13,0)}
\\&\qquad
+\big(1+1+1-1-1-1+1-1-1-1+1+1+1-1+1-1\big)\delta_{(14,0)}
\\&\qquad
+\big(i-i-i+i-i+i+i-i-i+i-i+i+i-i-i+i\big)\delta_{(15,0)}
\\&\qquad
+\big(-1+1-1+1-1-1+1+1+1-1-1+1-1+1-1+1\big)\delta_{(16,0)}
\\&\qquad
+\big(i+i-i-i+i+i+i-i+i-i-i-i+i+i-i-i\big)\delta_{(17,0)}\\&=17\delta_{(0 , 0)}\qedhere
\end{align*}
\end{proof}

\begin{rem} \label{remRealExtremtransform}
We note that if $\mu$ above is multiplied by $\delta_{-8}$, the resulting measure is
self-adjoint (that is, $(\delta_{-8}*\mu)\widetilde{\ \ } =\delta_{-8}*\mu$). Hence, the set 
$\{-8,\dots,8\}$ is the support of an extreme measure with real transform. See 
\propref{propExtrLessSingleton} for a related result.
\end{rem}
\newpage
%%%%%%%%%%%%%%%%%%%%%%%%%%%%%%%%%%%%%%%%%%%%%%%%
\section{Proofs of extremality for non-cyclic groups}\label{secnoncyclicgps}
%%%%%%%%%%%%%%%%%%%%%%%%%%%%%%%%%%%%%%%%%%%%%%%%



\subsection{\Zp23}
\begin{prop}
\cite{MR627683}
\label{prop5inZp2.Zp2.Zp2} The 5-element set 
$\{(0, 0, 0), (0, 0, 1), (0, 1, 0), (0, 1, 1), (1, 0, 0) \}$ 
is extreme in \Zp{2}{3}.
\end{prop}

 
\begin{proof}
Let  $\mu =  \delta_{(0,0,0)}-\delta_{(0,0,1)}-\delta_{(0,1,0)}-\delta_{(0,1,1)}+i\delta_{(1,0,0)}
$.  Then 
\begin{align*} 
\mu * \tilde \mu & = \big(1+1+1+1+1\big)\delta_{(0,0,0)}
\\&\qquad
+\big(-1-1+1+1\big)\delta_{(0,0,1)}
\\&\qquad
+\big(-1+1-1+1\big)\delta_{(0,1,0)}
\\&\qquad
+\big(-1+1+1-1\big)\delta_{(0,1,1)}
\\&\qquad
+\big(-i+i\big)\delta_{(1,0,0)}
\\&\qquad
+\big(i-i\big)\delta_{(1,0,1)}
\\&\qquad
+\big(i-i\big)\delta_{(1,1,0)}
\\&\qquad
+\big(i-i\big)\delta_{(1,1,1)}\\&=5\delta_{(0 , 0, 0)}\qedhere
\end{align*}
\end{proof}


\begin{prop}\label{prop6inZ2.X.Z2.X.Z2}
The 6 element set $\{(0, 0, 0),  (0, 0, 1),  (0, 1, 0),  (0, 1, 1),  (1, 0, 0),  (1, 0, 1) \}$ is extreme in $\Zp{2}{3}$. \newline \end{prop} 


 \begin{proof}
Let  $\mu =  \delta_{(0,0,0)}-i\delta_{(0,0,1)}+e^{7\pii/4}\delta_{(0,1,0)}+e^{3\pii/4}\delta_{(0,1,1)}+e^{\pii/4}\delta_{(1,0,0)}+e^{\pii/4}\delta_{(1,0,1)}
$.  Then 
\begin{align*} 
\mu * \tilde \mu & = \big(1+1+1+1+1+1\big)\delta_{(0,0,0)}
\\&\qquad
+\big(i-i-1-1+1+1\big)\delta_{(0,0,1)}
\\&\qquad
+\big(e^{\pii/4}+e^{3\pii/4}+e^{7\pii/4}+e^{5\pii/4}\big)\delta_{(0,1,0)}
\\&\qquad
+\big(e^{5\pii/4}+e^{7\pii/4}+e^{\pii/4}+e^{3\pii/4}\big)\delta_{(0,1,1)}
\\&\qquad
+\big(e^{7\pii/4}+e^{5\pii/4}+e^{\pii/4}+e^{3\pii/4}\big)\delta_{(1,0,0)}
\\&\qquad
+\big(e^{7\pii/4}+e^{5\pii/4}+e^{3\pii/4}+e^{\pii/4}\big)\delta_{(1,0,1)}
\\&\qquad
+\big(-i+i+i-i\big)\delta_{(1,1,0)}
\\&\qquad
+\big(-i+i-i+i\big)\delta_{(1,1,1)}\\&=6\delta_{(0 , 0, 0)}\qedhere
\end{align*}
\end{proof}




For \Zp23, a seven-element set has PSC   at most 2.777127870 unlike  seven-element subsets of \Zp8.



\subsection{\Zp32}

\begin{prop}\label{prop7inZp3.Zp3}
Each of the two 7-element sets below is extreme in \Zp{ 3}{2}.
\begin{enumerate}
\item  $\{(0, 0),  (0, 1),  (0, 2),  (1, 0),  (1, 1),  (2, 0),  (2, 2) \}$ 
\item $\{(0, 0),  (0, 1),  (0, 2),  (1, 0),  (1, 1),  (1, 2),  (2, 0) \}$
\end{enumerate} 

\end{prop}

\begin{proof}
(1). Let $\mu=\delta_{(0,0)}+e^{5\pii/3}\delta_{(0,1)}+e^{2\pii/3}\delta_{(0,2)}+e^{5\pii/3}\delta_{(1,0)}+\delta_{(1,1)}+e^{2\pii/3}\delta_{(2,0)}+e^{\pii/3}\delta_{(2,2)}
$.  Then 
\begin{align*} 
\mu * \tilde \mu & = \big(1+1+1+1+1+1+1\big)\delta_{(0,0)}
\\&
\quad+\big(e^{4\pii/3}+e^{5\pii/3}+e^{\pii/3}+e^{\pii/3}\big)\delta_{(0,1)}
\\&
\quad+\big(e^{\pii/3}+e^{2\pii/3}+e^{5\pii/3}+e^{5\pii/3}\big)\delta_{(0,2)}
\\&
\quad+\big(e^{4\pii/3}+e^{\pii/3}+e^{5\pii/3}+e^{\pii/3}-\big)\delta_{(1,0)}
\\&
\quad+\big(e^{5\pii/3}+e^{\pii/3}\big)\delta_{(1,1)}
\\&
\quad+\big(e^{4\pii/3}+e^{4\pii/3}+e^{2\pii/3}+e^{2\pii/3}\big)\delta_{(1,2)}
\\&
\quad+\big(e^{\pii/3}+e^{5\pii/3}+e^{2\pii/3}+e^{5\pii/3}\big)\delta_{(2,0)}
\\&
\quad+\big(e^{2\pii/3}+e^{4\pii/3}+e^{4\pii/3}+e^{2\pii/3}\big)\delta_{(2,1)}
\\&
\quad+\big(e^{5\pii/3}+e^{\pii/3}\big)\delta_{(2,2)}\\&=7\delta_{(0 , 0)}.
\end{align*}

(2). Let
$\mu=\delta_{(0,0)}+e^{4\pii/3}\delta_{(0,1)}+e^{4\pii/3}\delta_{(0,2)}+e^{5\pii/3}\delta_{(1,0)}+e^{\pii/3}\delta_{(1,1)}+e^{\pii/3}\delta_{(1,2)}+\delta_{(2,0)}
$.  Then 
\begin{align*} 
\mu * \tilde \mu & = \big(1+1+1+1+1+1+1\big)\delta_{(0,0)}
\\&
\quad+\big(e^{2\pii/3}+e^{4\pii/3}+1+e^{4\pii/3}+e^{2\pii/3}+1\big)\delta_{(0,1)}
\\&
\quad+\big(e^{2\pii/3}+1+e^{4\pii/3}+e^{4\pii/3}+1+e^{2\pii/3}\big)\delta_{(0,2)}
\\&
\quad+\big(1+e^{5\pii/3}-1-1+e^{\pii/3}\big)\delta_{(1,0)}
\\&
\quad+\big(e^{4\pii/3}+e^{\pii/3}+e^{\pii/3}-1+e^{5\pii/3}\big)\delta_{(1,1)}
\\&
\quad+\big(e^{4\pii/3}+e^{\pii/3}-1+e^{\pii/3}+e^{5\pii/3}\big)\delta_{(1,2)}
\\&
\quad+\big(e^{\pii/3}-1-1+e^{5\pii/3}+1\big)\delta_{(2,0)}
\\&
\quad+\big(e^{5\pii/3}+e^{5\pii/3}-1+e^{\pii/3}+e^{2\pii/3}\big)\delta_{(2,1)}
\\&
\quad+\big(e^{5\pii/3}-1+e^{5\pii/3}+e^{\pii/3}+e^{2\pii/3}\big)\delta_{(2,2)}\\&=7\delta_{(0 , 0)}.
\end{align*}
 \qedhere
\end{proof}
 

\begin{prop}\label{prop7inZ3.X.Z3}
The 7 element set $\{(0, 0),  (0, 1),  (0, 2),  (1, 0),  (1, 1),  (1, 2),  (2, 0) \}$ is extreme in $\Zp{3}{2}$. \newline \end{prop} 


 
\begin{proof}
Let  $\mu =  \delta_{(0,0)}+e^{4\pii/3}\delta_{(0,1)}+e^{4\pii/3}\delta_{(0,2)}+e^{\pii/3}\delta_{(1,0)}-\delta_{(1,1)}-\delta_{(1,2)}+e^{4\pii/3}\delta_{(2,0)}
$.  Then 
\begin{align*} 
\mu * \tilde \mu & = \big(1+1+1+1+1+1+1\big)\delta_{(0,0)}
\\&\qquad
+\big(e^{2\pii/3}+e^{4\pii/3}+1+e^{4\pii/3}+e^{2\pii/3}+1\big)\delta_{(0,1)}
\\&\qquad
+\big(e^{2\pii/3}+1+e^{4\pii/3}+e^{4\pii/3}+1+e^{2\pii/3}\big)\delta_{(0,2)}
\\&\qquad
+\big(e^{2\pii/3}+e^{\pii/3}+e^{5\pii/3}+e^{5\pii/3}-1\big)\delta_{(1,0)}
\\&\qquad
+\big(1-1-1+e^{5\pii/3}+e^{\pii/3}\big)\delta_{(1,1)}
\\&\qquad
+\big(1-1+e^{5\pii/3}-1+e^{\pii/3}\big)\delta_{(1,2)}
\\&\qquad
+\big(e^{5\pii/3}+e^{\pii/3}+e^{\pii/3}-1+e^{4\pii/3}\big)\delta_{(2,0)}
\\&\qquad
+\big(-1-1+e^{\pii/3}+e^{5\pii/3}+1\big)\delta_{(2,1)}
\\&\qquad
+\big(-1+e^{\pii/3}-1+e^{5\pii/3}+1\big)\delta_{(2,2)}\\&=7\delta_{(0 , 0)}\qedhere
\end{align*}
\end{proof}


% \subsection{\Zp33}

 
\subsection{\Zp{4}{2}.}

The following is a variant of 
\cite[3.3 (i) ]{MR627683}.

\begin{prop}\label{prop5inZ4.X.Z4}
Each of the following  5-element sets  is extreme in $  \Zp{4}{2}$. 
\begin{enumerate}
\item $\{(0, 0),  (0, 1),  (0, 2),  (0, 3),  (2, 0) \}$
\item $\{(0, 0),  (0, 1),  (0, 2),  (2, 0),  (2, 2) \}$
\end{enumerate}
 \end{prop} 


 
\begin{proof} (1).
Let  $\mu =  \delta_{(0,0)}-i\delta_{(0,1)}+\delta_{(0,2)}+i\delta_{(0,3)}+i\delta_{(2,0)}
$.  Then 
\begin{align*} 
\mu * \tilde \mu & = \big(1+1+1+1+1\big)\delta_{(0,0)}
%\\&\qquad
+\big(-i-i+i+i\big)\delta_{(0,1)}
\\&\qquad
+\big(1-1+1-1\big)\delta_{(0,2)}
\\&\qquad
+\big(i-i-i+i\big)\delta_{(0,3)}
%\\&\qquad
+\big(-i+i\big)\delta_{(2,0)}
\\&\qquad
+\big(-1+1\big)\delta_{(2,1)}
%\\&\qquad
+\big(-i+i\big)\delta_{(2,2)}
%\\&\qquad
+\big(1-1\big)\delta_{(2,3)}\\&=5\delta_{(0 , 0)}.
\end{align*}
\

(2).
Let  $\mu =  \delta_{(0,0)}-\delta_{(0,1)}-\delta_{(0,2)}+i\delta_{(2,0)}+i\delta_{(2,2)}
$.  Then 
\begin{align*} 
\mu * \tilde \mu & = \big(1+1+1+1+1\big)\delta_{(0,0)}
%\\&\qquad
+\big(-1+1\big)\delta_{(0,1)}
\\&\qquad
+\big(-1-1+1+1\big)\delta_{(0,2)}
%\\&\qquad
+\big(-1+1\big)\delta_{(0,3)}
\\&\qquad
+\big(-i+i+i-i\big)\delta_{(2,0)}
%\\&\qquad
+\big(i-i\big)\delta_{(2,1)}
\\&\qquad
+\big(-i+i-i+i\big)\delta_{(2,2)}
%\\&\qquad
+\big(i-i\big)\delta_{(2,3)}\\&=5\delta_{(0 , 0)}\qedhere
\end{align*}
\end{proof}




\begin{prop}\label{prop6inZ4.X.Z4}
Each of the 6-element sets 
\begin{enumerate}
\item $\{(0, 0),  (0, 1),  (0, 2),  (0, 3),  (1, 0),  (1, 2) \}$
\item $\{(0, 0),  (0, 1),  (0, 2),  (1, 0),  (1, 1),  (1, 2) \}$
\end{enumerate}
is extreme in $\Zp2{} \times \Zp{4}{}$.  
\end{prop} 

\begin{rem}
The 2018/10/30 version of my program (findNonCyclic) did not find those two sets equivalent. A better program might, however.
\end{rem}
 

 
\begin{proof} (1). 
Let  $\mu =  \delta_{(0,0)}+e^{7\pii/4}\delta_{(0,1)}+i\delta_{(0,2)}+e^{3\pii/4}\delta_{(0,3)}+e^{3\pii/4}\delta_{(1,0)}+e^{3\pii/4}\delta_{(1,2)}
$.  Then 
\begin{align*} 
\mu * \tilde \mu & = \big(1+1+1+1+1+1\big)\delta_{(0,0)}
\\&\qquad
+\big(e^{5\pii/4}+e^{7\pii/4}+e^{3\pii/4}+e^{\pii/4}\big)\delta_{(0,1)}
\\&\qquad
+\big(-i-1+i-1+1+1\big)\delta_{(0,2)}
\\&\qquad
+\big(e^{\pii/4}+e^{5\pii/4}+e^{7\pii/4}+e^{3\pii/4}\big)\delta_{(0,3)}
\\&\qquad
+\big(e^{5\pii/4}+e^{7\pii/4}+e^{3\pii/4}+e^{\pii/4}\big)\delta_{(1,0)}
\\&\qquad
+\big(-1+1+1-1\big)\delta_{(1,1)}
\\&\qquad
+\big(e^{5\pii/4}+e^{7\pii/4}+e^{\pii/4}+e^{3\pii/4}\big)\delta_{(1,2)}
\\&\qquad
+\big(-1+1-1+1\big)\delta_{(1,3)}\\&=6\delta_{(0 , 0)}\qedhere
\end{align*}


(2). 
Let  $\mu =  \delta_{(0,0)}+e^{23\pii/12}\delta_{(0,1)}-i\delta_{(0,2)}+e^{\pii/3}\delta_{(1,0)}+e^{11\pii/12}\delta_{(1,1)}+e^{11\pii/6}\delta_{(1,2)}
$.  Then 
\begin{align*} 
\mu * \tilde \mu & = \big(1+1+1+1+1+1\big)\delta_{(0,0)}
\\&\qquad
+\big(e^{23\pii/12}+e^{19\pii/12}+e^{7\pii/12}+e^{11\pii/12}\big)\delta_{(0,1)}
\\&\qquad
+\big(i-i+i-i\big)\delta_{(0,2)}
\\&\qquad
+\big(e^{\pii/12}+e^{5\pii/12}+e^{17\pii/12}+e^{13\pii/12}\big)\delta_{(0,3)}
\\&\qquad
+\big(e^{5\pii/3}-1+e^{5\pii/3}+e^{\pii/3}-1+e^{\pii/3}\big)\delta_{(1,0)}
\\&\qquad
+\big(e^{19\pii/12}+e^{7\pii/12}+e^{11\pii/12}+e^{23\pii/12}\big)\delta_{(1,1)}
\\&\qquad
+\big(e^{\pii/6}+e^{7\pii/6}+e^{5\pii/6}+e^{11\pii/6}\big)\delta_{(1,2)}
\\&\qquad
+\big(e^{13\pii/12}+e^{\pii/12}+e^{5\pii/12}+e^{17\pii/12}\big)\delta_{(1,3)}\\&=6\delta_{(0 , 0)}\qedhere
\end{align*}
\end{proof}

%\section{\Zp k n \ for $k\ge 4$ and $n\ge 2$ (aspirational)}
\subsection{\Zp43}
These groups have 25 or more elements, too many to find equivalence classes for sets with 6 or more elements in a
reasonable time.

%\section{More complicated products}

\subsection{$\Zp{2}{}\times\Zp4{}$}


\begin{prop}\label{prop4inZ2.X.Z4}
The 4 element set $\{(0, 0),  (0, 1),  (1, 3),  (1, 0) \}$ is extreme in $\Zp{2}{} \times \Zp{4}{}$. \newline \end{prop} 

 
\begin{proof}
Let  $\mu =  \delta_{(0,0)}+i\delta_{(0,1)}-\delta_{(1,3)}+i\delta_{(1,0)}
$.  Then 
\begin{align*} 
\mu * \tilde \mu & = \big(1+1+1+1\big)\delta_{(0,0)}
\\&\qquad
+\big(i-i\big)\delta_{(0,1)}
\\&\qquad
+\big(-i+i\big)\delta_{(0,3)}
\\&\qquad
+\big(-i+i\big)\delta_{(1,0)}
\\&\qquad
+\big(-1+1\big)\delta_{(1,1)}
\\&\qquad
+\big(-i+i\big)\delta_{(1,2)}
\\&\qquad
+\big(-1+1\big)\delta_{(1,3)}\\&=4\delta_{(0 , 0)}\qedhere
\end{align*}
\end{proof}



\begin{prop}\label{prop5inZ2.X.Z4}
Each of the following 5-element sets is extreme in \Zp2{}\Times\Zp4{}.
\begin{enumerate}

%\item $\{(0, 0),  (0, 1),  (0, 2),  (1, 1),  (1, 3) \}$ 
\item $\{(0, 0),  (0, 1),  (0, 2),  (1, 0),  (1, 2) \}$
\item $\{(0, 0),  (0, 1),  (0, 2),  (0, 3),  (1, 0) \}$ 
\end{enumerate}
 \end{prop} 
 
\begin{proof}
(1). 
\begin{comment}
Let  $\mu =  \delta_{(0,0)}+\delta_{(0,1)}-\delta_{(0,2)}+i\delta_{(1,1)}+i\delta_{(1,3)}
$.  Then 
\begin{align*} 
\mu * \tilde \mu & = \big(1+1+1+1+1\big)\delta_{(0,0)}
%\\&\qquad
+\big(1-1\big)\delta_{(0,1)}
%\\&\qquad
+\big(-1-1+1+1\big)\delta_{(0,2)}
\\&\qquad
+\big(1-1\big)\delta_{(0,3)}
%\\&\qquad
+\big(-i+i\big)\delta_{(1,0)}
%\\&\qquad
+\big(-i+i+i-i\big)\delta_{(1,1)}
\\&\qquad
+\big(-i+i\big)\delta_{(1,2)}
%\\&\qquad
+\big(-i+i-i+i\big)\delta_{(1,3)}\\&=5\delta_{(0 , 0)}.
\end{align*}

(2).
\end{comment}
Let  $\mu =  \delta_{(0,0)}-\delta_{(0,1)}-\delta_{(0,2)}+i\delta_{(1,0)}+i\delta_{(1,2)}
$.  Then 
\begin{align*} 
\mu * \tilde \mu & = \big(1+1+1+1+1\big)\delta_{(0,0)}
%\\&\qquad
+\big(-1+1\big)\delta_{(0,1)}
%\\&\qquad
+\big(-1-1+1+1\big)\delta_{(0,2)}
\\&\qquad
+\big(-1+1\big)\delta_{(0,3)}
%\\&\qquad
+\big(-i+i+i-i\big)\delta_{(1,0)}
%\\&\qquad
+\big(i-i\big)\delta_{(1,1)}
\\&\qquad
+\big(-i+i-i+i\big)\delta_{(1,2)}
%\\&\qquad
+\big(i-i\big)\delta_{(1,3)}\\&=5\delta_{(0 , 0)}.
\end{align*}
 
 (2). 
Let  $\mu =  \delta_{(0,0)}-i\delta_{(0,1)}+\delta_{(0,2)}+i\delta_{(0,3)}+i\delta_{(1,0)}
$.  Then 
\begin{align*} 
\mu * \tilde \mu & = \big(1+1+1+1+1\big)\delta_{(0,0)}
%\\&\qquad
+\big(-i-i+i+i\big)\delta_{(0,1)}
\\&\qquad
+\big(1-1+1-1\big)\delta_{(0,2)}
%\\&\qquad
+\big(i-i-i+i\big)\delta_{(0,3)}
%\\&\qquad
+\big(-i+i\big)\delta_{(1,0)}
\\&\qquad
+\big(-1+1\big)\delta_{(1,1)}
%\\&\qquad
+\big(-i+i\big)\delta_{(1,2)}
%\\&\qquad
+\big(1-1\big)\delta_{(1,3)}\\&=5\delta_{(0 , 0)}\qedhere
\end{align*}
\end{proof}



\begin{prop}\label{prop6inZ2.X.Z4}
Each of the following 6-element sets  
is  extreme in $\Zp{2}{} \times \Zp{4}{}$. 
\begin{enumerate}
\item$\{(0, 0),  (0, 1),  (0, 2),  (0, 3),  (1, 0),  (1, 2) \}$
\item  $\{(0, 0),  (0, 1),  (0, 2),  (1, 0),  (1, 1),  (1, 2) \}$
\end{enumerate}
 \end{prop} 


 
\begin{proof}
(1.)
Let  $\mu =  \delta_{(0,0)}+e^{7\pii/4}\delta_{(0,1)}+i\delta_{(0,2)}+e^{3\pii/4}\delta_{(0,3)}+e^{3\pii/4}\delta_{(1,0)}+e^{3\pii/4}\delta_{(1,2)}
$.  Then 
\begin{align*} 
\mu * \tilde \mu & = \big(1+1+1+1+1+1\big)\delta_{(0,0)}
\\&\qquad
+\big(e^{5\pii/4}+e^{7\pii/4}+e^{3\pii/4}+e^{\pii/4}\big)\delta_{(0,1)}
\\&\qquad
+\big(-i-1+i-1+1+1\big)\delta_{(0,2)}
\\&\qquad
+\big(e^{\pii/4}+e^{5\pii/4}+e^{7\pii/4}+e^{3\pii/4}\big)\delta_{(0,3)}
\\&\qquad
+\big(e^{5\pii/4}+e^{7\pii/4}+e^{3\pii/4}+e^{\pii/4}\big)\delta_{(1,0)}
%\\&\qquad
+\big(-1+1+1-1\big)\delta_{(1,1)}
\\&\qquad
+\big(e^{5\pii/4}+e^{7\pii/4}+e^{\pii/4}+e^{3\pii/4}\big)\delta_{(1,2)}
%\\&\qquad
+\big(-1+1-1+1\big)\delta_{(1,3)}\\&=6\delta_{(0 , 0)}.
\end{align*}

(2).
Let  $\mu =  \delta_{(0,0)}+e^{23\pii/12}\delta_{(0,1)}-i\delta_{(0,2)}+e^{\pii/3}\delta_{(1,0)}+e^{11\pii/12}\delta_{(1,1)}+e^{11\pii/6}\delta_{(1,2)}
$.  Then 
\begin{align*} 
\mu * \tilde \mu & = \big(1+1+1+1+1+1\big)\delta_{(0,0)}
\\&\qquad
+\big(e^{23\pii/12}+e^{19\pii/12}+e^{7\pii/12}+e^{11\pii/12}\big)\delta_{(0,1)}
\\&\qquad
+\big(i-i+i-i\big)\delta_{(0,2)}
\\&\qquad
+\big(e^{\pii/12}+e^{5\pii/12}+e^{17\pii/12}+e^{13\pii/12}\big)\delta_{(0,3)}
\\&\qquad
+\big(e^{5\pii/3}-1+e^{5\pii/3}+e^{\pii/3}-1+e^{\pii/3}\big)\delta_{(1,0)}
\\&\qquad
+\big(e^{19\pii/12}+e^{7\pii/12}+e^{11\pii/12}+e^{23\pii/12}\big)\delta_{(1,1)}
\\&\qquad
+\big(e^{\pii/6}+e^{7\pii/6}+e^{5\pii/6}+e^{11\pii/6}\big)\delta_{(1,2)}
\\&\qquad
+\big(e^{13\pii/12}+e^{\pii/12}+e^{5\pii/12}+e^{17\pii/12}\big)\delta_{(1,3)}\\&=6\delta_{(0 , 0)}\qedhere
\end{align*}
\end{proof}





\begin{comment} % prop6inZ2.Z4.Z4 contains prop6inZ23
\subsection{$\Zp{2}{}\times\Zp{4}{2}$}

\marginpar{\tiny How is this Prop related to  \propref{prop6inZ4.X.Z4} and 
\propref{prop6inZp2.Zp4.Zp4}?}

\begin{prop}\label{prop6inZp2.Zp4.Zp4}
Each of the following 6-element sets is extreme in $\Zp{2}{} \times \Zp{4}{2}$.
\begin{enumerate}
\item  $\{(0, 0, 0),  (0, 0, 1),  (0, 0, 2),  (0, 0, 3),  (1, 0, 0),  (1, 0, 2) \}$ %1
\item  $\{(0, 0, 0),  (0, 0, 1),  (0, 0, 2),  (0, 0, 3),  (1, 1, 0),  (1, 3, 2) \}$ %2
\item  $\{(0, 0, 0),  (0, 0, 1),  (0, 0, 2),  (0, 0, 3),  (1, 2, 0),  (1, 2, 2) \}$ %3
\item  $\{(0, 0, 0),  (0, 0, 1),  (0, 0, 2),  (0, 0, 3),  (1, 1, 0),  (1, 3, 0) \}$ %4
\end{enumerate}
%They are non-equivalent.
\end{prop} 


 
\begin{proof}
(1). 
Let  $\mu =  \delta_{(0,0,0)}+e^{7\pii/4}\delta_{(0,0,1)}+i\delta_{(0,0,2)}+e^{3\pii/4}\delta_{(0,0,3)}+e^{3\pii/4}\delta_{(1,0,0)}+e^{3\pii/4}\delta_{(1,0,2)}
$.  Then 
\begin{align*} 
\mu * \tilde \mu & = \big(1+1+1+1+1+1\big)\delta_{(0,0,0)}\\&\qquad
+\big(e^{5\pii/4}+e^{7\pii/4}+e^{3\pii/4}+e^{\pii/4}\big)\delta_{(0,0,1)}
\\&\qquad
+\big(-i-1+i-1+1+1\big)\delta_{(0,0,2)}
\\&\qquad
+\big(e^{\pii/4}+e^{5\pii/4}+e^{7\pii/4}+e^{3\pii/4}\big)\delta_{(0,0,3)}
\\&\qquad
+\big(e^{5\pii/4}+e^{7\pii/4}+e^{3\pii/4}+e^{\pii/4}\big)\delta_{(1,0,0)}
\\&\qquad
+\big(-1+1+1-1\big)\delta_{(1,0,1)}
\\&\qquad
+\big(e^{5\pii/4}+e^{7\pii/4}+e^{\pii/4}+e^{3\pii/4}\big)\delta_{(1,0,2)}
\\&\qquad
+\big(-1+1-1+1\big)\delta_{(1,0,3)}\\&=6\delta_{(0 , 0, 0)}.
\end{align*}

(2). 
Let  $\mu =  \delta_{(0,0,0)}+\delta_{(0,0,1)}+\delta_{(0,0,2)}-\delta_{(0,0,3)}+e^{\pii/4}\delta_{(1,1,0)}+e^{3\pii/4}\delta_{(1,3,2)}
$.  Then 
\begin{align*} 
\mu * \tilde \mu & = \big(1+1+1+1+1+1\big)\delta_{(0,0,0)}
%\\&\qquad
+\big(-1+1+1-1\big)\delta_{(0,0,1)}
\\&\qquad
+\big(1-1+1-1\big)\delta_{(0,0,2)}
%\\&\qquad
+\big(1+1-1-1\big)\delta_{(0,0,3)}
\\&\qquad
+\big(-i+i\big)\delta_{(0,2,2)}
%\\&\qquad
+\big(e^{5\pii/4}+e^{\pii/4}\big)\delta_{(1,1,0)}
\\&\qquad
+\big(e^{\pii/4}+e^{5\pii/4}\big)\delta_{(1,1,1)}
%\\&\qquad
+\big(e^{5\pii/4}+e^{\pii/4}\big)\delta_{(1,1,2)}
\\&\qquad
+\big(e^{5\pii/4}+e^{\pii/4}\big)\delta_{(1,1,3)}
%\\&\qquad
+\big(e^{7\pii/4}+e^{3\pii/4}\big)\delta_{(1,3,0)}
\\&\qquad
+\big(e^{7\pii/4}+e^{3\pii/4}\big)\delta_{(1,3,1)}
%\\&\qquad
+\big(e^{7\pii/4}+e^{3\pii/4}\big)\delta_{(1,3,2)}
\\&\qquad
+\big(e^{3\pii/4}+e^{7\pii/4}\big)\delta_{(1,3,3)}\\&=6\delta_{(0 , 0, 0)}.
\end{align*}
(3). 
Let  $\mu =  \delta_{(0,0,0)}+e^{7\pii/4}\delta_{(0,0,1)}+i\delta_{(0,0,2)}+e^{3\pii/4}\delta_{(0,0,3)}+e^{3\pii/4}\delta_{(1,2,0)}+e^{3\pii/4}\delta_{(1,2,2)}
$.  Then 
\begin{align*} 
\mu * \tilde \mu & = \big(1+1+1+1+1+1\big)\delta_{(0,0,0)}
\\&\qquad
+\big(e^{5\pii/4}+e^{7\pii/4}+e^{3\pii/4}+e^{\pii/4}\big)\delta_{(0,0,1)}
\\&\qquad
+\big(-i-1+i-1+1+1\big)\delta_{(0,0,2)}
\\&\qquad
+\big(e^{\pii/4}+e^{5\pii/4}+e^{7\pii/4}+e^{3\pii/4}\big)\delta_{(0,0,3)}
\\&\qquad
+\big(e^{5\pii/4}+e^{7\pii/4}+e^{3\pii/4}+e^{\pii/4}\big)\delta_{(1,2,0)}
\\&\qquad
+\big(-1+1+1-1\big)\delta_{(1,2,1)}
\\&\qquad
+\big(e^{5\pii/4}+e^{7\pii/4}+e^{\pii/4}+e^{3\pii/4}\big)\delta_{(1,2,2)}
\\&\qquad
+\big(-1+1-1+1\big)\delta_{(1,2,3)}\\&=6\delta_{(0 , 0, 0)}.
\end{align*}

 
(4).
Let  $\mu =  \delta_{(0,0,0)}-\delta_{(0,0,1)}-\delta_{(0,0,2)}-\delta_{(0,0,3)}+e^{3\pii/4}\delta_{(1,1,0)}+e^{3\pii/4}\delta_{(1,3,0)}
$.  Then 
\begin{align*} 
\mu * \tilde \mu & = \big(1+1+1+1+1+1\big)\delta_{(0,0,0)}
%\\&\qquad
+\big(-1-1+1+1\big)\delta_{(0,0,1)}
\\&\qquad
+\big(-1+1-1+1\big)\delta_{(0,0,2)}
%\\&\qquad
+\big(-1+1+1-1\big)\delta_{(0,0,3)}
\\&\qquad
+\big(1+1\big)\delta_{(0,2,0)}
%\\&\qquad
+\big(e^{5\pii/4}+e^{3\pii/4}\big)\delta_{(1,1,0)}
\\&\qquad
+\big(e^{\pii/4}+e^{7\pii/4}\big)\delta_{(1,1,1)}
%\\&\qquad
+\big(e^{\pii/4}+e^{7\pii/4}\big)\delta_{(1,1,2)}
\\&\qquad
+\big(e^{\pii/4}+e^{7\pii/4}\big)\delta_{(1,1,3)}
%\\&\qquad
+\big(e^{5\pii/4}+e^{3\pii/4}\big)\delta_{(1,3,0)}
\\&\qquad
+\big(e^{\pii/4}+e^{7\pii/4}\big)\delta_{(1,3,1)}
%\\&\qquad
+\big(e^{\pii/4}+e^{7\pii/4}\big)\delta_{(1,3,2)}
\\&\qquad
+\big(e^{\pii/4}+e^{7\pii/4}\big)\delta_{(1,3,3)}\\&=6\delta_{(0 , 0, 0)}.
\qedhere
\end{align*}

\end{proof}
\end{comment}

\subsection{\Zp{2}{2}\Times\Zp{3}{}}

\begin{prop}\label{prop8inZ2.X.Z2.X.Z3}
The 8 element set 
\[
\{(0, 0, 0),  (0, 0, 1),  (0, 1, 0),  (0, 1, 1),  (1, 0, 0),  (1, 0, 1),  (1, 1, 0),  (1, 1, 1) \}
\] 
is extreme in $\Zp{2}{2}\Zp{3}{}$. \newline \end{prop} 


 
\begin{proof}
Let  $\mu =  \delta_{(0,0,0)}+e^{7\pii/15}\delta_{(0,0,1)}-i\delta_{(0,1,0)}+e^{29\pii/30}\delta_{(0,1,1)}+i\delta_{(1,0,0)}+e^{29\pii/30}\delta_{(1,0,1)}-\delta_{(1,1,0)}+e^{7\pii/15}\delta_{(1,1,1)}
$.  Then 
\begin{align*} 
\mu * \tilde \mu & = \big(1+1+1+1+1+1+1+1\big)\delta_{(0,0,0)}
\\&\qquad
+\big(e^{7\pii/15}+e^{22\pii/15}+e^{7\pii/15}+e^{22\pii/15}\big)\delta_{(0,0,1)}
\\&\qquad
+\big(e^{23\pii/15}+e^{8\pii/15}+e^{23\pii/15}+e^{8\pii/15}\big)\delta_{(0,0,2)}
\\&\qquad
+\big(i-i-i+i-i+i+i-i\big)\delta_{(0,1,0)}
\\&\qquad
+\big(e^{29\pii/30}+e^{29\pii/30}+e^{59\pii/30}+e^{59\pii/30}\big)\delta_{(0,1,1)}
\\&\qquad
+\big(e^{31\pii/30}+e^{31\pii/30}+e^{\pii/30}+e^{\pii/30}\big)\delta_{(0,1,2)}
\\&\qquad
+\big(-i-i+i+i+i+i-i-i\big)\delta_{(1,0,0)}
\\&\qquad
+\big(e^{59\pii/30}+e^{59\pii/30}+e^{29\pii/30}+e^{29\pii/30}\big)\delta_{(1,0,1)}
\\&\qquad
+\big(e^{31\pii/30}+e^{31\pii/30}+e^{\pii/30}+e^{\pii/30}\big)\delta_{(1,0,2)}
\\&\qquad
+\big(-1+1-1+1-1+1-1+1\big)\delta_{(1,1,0)}
\\&\qquad
+\big(e^{22\pii/15}+e^{7\pii/15}+e^{22\pii/15}+e^{7\pii/15}\big)\delta_{(1,1,1)}
\\&\qquad
+\big(e^{23\pii/15}+e^{8\pii/15}+e^{23\pii/15}+e^{8\pii/15}\big)\delta_{(1,1,2)}\\&=8\delta_{(0 , 0, 0)}\qedhere
\end{align*}
\end{proof}

\begin{comment}
\subsection{Perhaps to do}
\begin{verbatim}

As results are obtained for these groups,
 these sections will have material added to them and 
 be moved into the regular text.
\subsection{$\Zp{2}{}\times\Zp{5}{2}$}
\subsection{$\Zp22 \times\Zp3{}{}$}
\subsection{$\Zp22 \times\Zp4{}{}$}
\subsection{$\Zp22 \times\Zp5{}{}$}
\subsection{$\Zp22 \times\Zp7{}{}$}
\subsection{$\Zp{2}{}\times\Zp{3}{2}$}
\subsection{$\Zp{3}{2}\times\Zp{4}{}$}
\subsection{$\Zp{4}{2}\times\Zp{3}{}$}
\end{verbatim}
\end{comment} 


%\input{bestiary-nonextrsets.tex}


%\input{bestiary-boring.tex}
%\input{bestiary-programming.tex}


\bibliographystyle{amsplain} %{plainnat}
\bibliography{ccg}


\newpage
\section{Disclaimers} 
{\Small
\subsection*{\ \ A request}  Suggestions for things to look at, programming ideas, and  corrections to the text are invited.



\subsection*{\ \  Conjectures} I have mostly used ``conjecture'' rather than ``problem'' or ``question,'' in the hope
that ``conjecture'' will pique competitive instincts in a way that the other terms may not. 

\subsection*{\ \ To do - theory}
For  each   conjecture, show that it is correct (or give counter example(s) to it). In some cases this will
surely involve rephrasing of the conjecture.
\smallbreak

Criteria for ``exceptional'' extreme sets in non-cyclic groups need formulating. 

\subsection*{\ \ To do - computation} 
There is much computer work to be done for larger sets in larger  
groups, where the results are limited by the slowness of the computer and my
meagre programming skills.  

 
\smallbreak
 A related project is to produce lists of extremal measures for each extreme set, in the hope of revealing  a
pattern (or patterns) that would enable solution of one or more of the conjectures. 
This would take relatively minor adaptations of the present programs.

\smallbreak
 Another project is to compile a list of sets with (close to) integer PSCs and then verify that the PSCs of
those sets are indeed integers.

\subsection*{\ \ Too much detail}
I don't expect anyone to read everything, 
  but the data is here if for those who want (excessively) repetitive exercises with answers.
  
  Also, some of what's here duplicates, partially, material of 
  the two papers in the bibliography of which I am an author or co-author. 
  
\subsection*{\ \ Are sets actually equivalent?}  Yes, for cyclic groups.  Perhaps not, for product groups: implementing 
in product groups a correct method for  finding true equivalence classes is 
 complicated and may well have escaped me; see, e.g., 
\cite{MR2363058}. Any help on this point would be much appreciated.

 
\section{\ \ The computer programs} I will be happy to send my computer programs to anyone who wants them. 


 

\bibliographystyle{amsplain} %{plainnat}
\bibliography{ccg}


\end{document}
%\bibliographystyle{amsplain}\bibliography{cm5558}
%%%%%%%%%%%%%%%%%%%%%%%%%%%%%%%%%%%%%%%%%%%%%%%%%%%%
 \newpage
\end{document}
